\documentclass[11pt]{article}

\usepackage[margin=1in]{geometry}
\usepackage{amsmath,amssymb,amsthm,mathtools,mathrsfs}
\usepackage{microtype}
\usepackage[hidelinks]{hyperref}
\allowdisplaybreaks

\newtheorem{theorem}{Theorem}
\newtheorem{lemma}[theorem]{Lemma}
\newtheorem{proposition}[theorem]{Proposition}
\newtheorem{corollary}[theorem]{Corollary}
\newtheorem{claim}[theorem]{Claim}
\theoremstyle{definition}
\newtheorem{definition}[theorem]{Definition}
\theoremstyle{remark}
\newtheorem{remark}[theorem]{Remark}

\newcommand{\I}{[0,1]}
\newcommand{\Fix}{\operatorname{Fix}}
\newcommand{\Crit}{\operatorname{Turn}}
\newcommand{\id}{\operatorname{id}}
\newcommand{\cl}{\operatorname{cl}}
\newcommand{\M}{\mathcal M_0}
\newcommand{\T}{\mathcal T}
\newcommand{\Lap}{\mathscr L}
\newcommand{\Var}{\operatorname{Var}}

\title{Common Fixed Points for Commuting\\
Finite-Piece Piecewise-Monotone Interval Maps}
\author{}
\date{}

\begin{document}
\maketitle

\begin{abstract}
Let $P,Q:[0,1]\to[0,1]$ be commuting continuous maps, each monotone on
finitely many interval pieces.  Monotonicity need not be strict, so constant
branches and plateau intervals are allowed.  We prove that $P$ and $Q$ have
a common fixed point.  The proof constructs a common monotone quotient on
which the two maps are finite-piece piecewise linear.  The required
piecewise-linear fixed-point proposition is proved directly by combining a
minimal-counterexample reduction, one-sided orbit-circle rigidity, exactness,
a common constant-slope coordinate, substitution dynamics, and an affine
circle argument.
\end{abstract}

\tableofcontents

\section{Introduction}

A continuous map $T:\I\to\I$ is called \emph{finite-piece piecewise
monotone} if there is a finite partition
\[
 0=x_0<x_1<\cdots <x_m=1
\]
such that $T$ is nondecreasing or nonincreasing on every
$[x_{i-1},x_i]$.  Constant pieces and nontrivial plateau intervals are
allowed.  The class is closed under composition.  To verify this last
assertion without a strict-monotonicity assumption, let \(U,V\) have finite
monotonicity partitions.  On a monotonicity cell \(J\) of \(V\), the inverse
image of each partition point of \(U\) is empty, a point, or a compact
interval, because \(V|_J\) is continuous and monotone.  Add the endpoints of
all these inverse images to the partition of \(V\).  On each cell of the
resulting finite refinement, \(V\) has image in one monotonicity cell of \(U\)
(or is constant at a partition point).  The composite \(U\circ V\) is
therefore monotone on that cell.  This proves the asserted closure, including
all plateau cases.

\begin{theorem}\label{thm:main}
Let $P,Q:\I\to\I$ be commuting continuous finite-piece
piecewise-monotone maps.  Then
\[
 \Fix(P)\cap\Fix(Q)\ne\varnothing.
\]
\end{theorem}

The corresponding statement for arbitrary continuous interval maps is
false by examples of Boyce and Huneke \cite{Boyce,Huneke}.  The proof of
Theorem~\ref{thm:main} has two necessary components.  First, we prove the
piecewise-linear core proposition below.  Second, using positive entropy and
a branch-transfer operator, we construct one nonatomic conformal measure for
both maps.  Its distribution function collapses all plateaus and descends
the pair to commuting finite-piece piecewise-linear maps.  A fixed point in
the quotient gives a common invariant fiber, and minimality finishes the
argument.

\section{The piecewise-linear core}

A map is \emph{finite-piece piecewise linear} if it is continuous and affine
on every interval of some finite partition; constant affine pieces are
allowed.
Throughout the proof, products of maps are compositions read from right to
left:
\[
 UV=U\circ V,\qquad (UV)(x)=U(V(x)).
\]

\begin{proposition}[Piecewise-linear core]\label{prop:PL-core}
Every pair of commuting continuous finite-piece piecewise-linear selfmaps of
$\I$ has a common fixed point.
\end{proposition}

\par\noindent\textit{Proof of Proposition~\ref{prop:PL-core}.}\quad
Suppose, toward a contradiction, that $P,Q$ commute and have no common fixed
point.

\subsection{Piecewise-linear finiteness and safe quotients}

\begin{lemma}[Piecewise-linear finiteness]\label{lem:pl-finiteness}
The class of continuous finite-piece piecewise-linear interval maps has the
following properties.
\begin{enumerate}
\item It is closed under composition, restriction to a compact interval,
and affine changes of coordinates.
\item For every map $T$ in the class, $\Fix(T)$ is a finite union of points
and compact intervals.
\item If $T$ has no flat interval, then it is finite-to-one and has finitely
many strict monotonicity laps.
\item On every affine branch, the derivative is constant and therefore has
bounded variation.
\item Suppose finitely many pairwise disjoint nondegenerate compact intervals
are collapsed by the canonical finite piecewise-affine nondecreasing
quotient, and a finite family of maps respects the quotient fibres.  The
descended maps are again continuous finite-piece piecewise linear.
\end{enumerate}
\end{lemma}

\begin{proof}
Composition, restriction, and affine conjugacy preserve finite piecewise
linearity after a finite refinement of the relevant partitions.  On each
affine branch, $T-\id$ is affine; hence its zero set is empty, a point, or
the whole branch.  This proves (2).  If $T$ has no flat interval, every
nondegenerate affine branch is strict, so only finitely many monotonicity
laps occur and every fiber meets each lap in at most one point.  This proves
(3), while (4) is immediate.

For (5), let $\pi$ be the quotient map.  On every uncollapsed source gap,
$\pi$ has an affine inverse.  After refining at the branch endpoints of the
original map and their images under $\pi$, the descended map is therefore a
composition of affine maps on each nondegenerate quotient interval.  On a
collapsed source interval, respect of fibres makes $\pi T$ constant.  The
point-quotient case gives a constant map.  Thus the descended map is
continuous and finite-piece piecewise linear.
\end{proof}

Here is the quotient convention used throughout.  Given pairwise disjoint
nondegenerate intervals $C_k=[a_k,b_k]$, $1\leq k\leq m$, put
\[
 d=1-\sum_{k=1}^m(b_k-a_k).
\]
If $d>0$, their \emph{canonical quotient} is the map $\pi:\I\to\I$ defined
by
\begin{equation}\label{eq:canonical-quotient}
 \pi(x)=\frac{x-\sum_{k=1}^m
              \operatorname{length}(C_k\cap[0,x])}{d}.
\end{equation}
It collapses precisely the intervals $C_k$ and is affine with slope $d^{-1}$
on every complementary gap.  If $d=0$, necessarily the collapsed family is
$\{\I\}$; the target is the one-point interval and $\pi$ is the constant
map.  A map $F$ \emph{respects the fibres of $\pi$} when
\begin{equation}\label{eq:respect-fibres}
 \pi(x)=\pi(y)\quad\Longrightarrow\quad
 \pi(F(x))=\pi(F(y)).
\end{equation}
Equivalently, $\pi F$ is constant on every quotient fibre.

Here and below, an interval \emph{invariant under \(T\)} is forward
invariant, \(T(J)\subseteq J\), unless equality is explicitly asserted.  Let
$\mathscr J$ be the family of nonempty compact intervals invariant under
both $P$ and $Q$, ordered by inclusion.  It is nonempty because it contains
$\I$.  If $\mathscr C\subseteq\mathscr J$ is a chain, every finite subfamily
of $\mathscr C$ has nonempty intersection: its intersection is simply its
least member.  Compactness of $\I$ and the finite-intersection property give
\(J_*=\bigcap_{J\in\mathscr C}J\neq\varnothing\).  This is a compact
interval, and for $x\in J_*$ and $J\in\mathscr C$ one has
$P(x),Q(x)\in J$; hence $P(J_*),Q(J_*)\subseteq J_*$.  Thus $J_*$ is a lower
bound, and Zorn's lemma gives an inclusion-minimal member.  Restrict to it
and rescale affinely to $\I$.  It is nondegenerate, since a singleton would
be a common fixed point.  Moreover
\begin{equation}\label{eq:onto}
 P(\I)=\I=Q(\I).
\end{equation}
Indeed, $P(\I)$ is a nonempty compact interval and
\[
 P(P(\I))\subseteq P(\I),\qquad
 Q(P(\I))=P(Q(\I))\subseteq P(\I).
\]
Minimality gives $P(\I)=\I$; interchanging $P,Q$ proves $Q(\I)=\I$.
Henceforth there is no proper nonempty compact interval invariant under both
maps.

\begin{proposition}[Finite respected quotient package]
\label{prop:quotient-package}
Let $\pi:\I\to I'$ be the canonical nondecreasing quotient which collapses
finitely many pairwise disjoint nondegenerate compact intervals and is
affine with positive slope on every complementary gap (with the one-point
case understood as above).  Let $\mathscr F$ be a finite commuting
family of maps in the finite-piece piecewise-linear class, and suppose every member of
$\mathscr F$ respects the fibres of $\pi$.  Then:
\begin{enumerate}
\item every $F\in\mathscr F$ descends to a unique continuous map
$\overline F:I'\to I'$ in the same piecewise-linear class; the descended maps commute,
and surjectivity descends;
\item for every map $T$ respecting the fibres,
\begin{equation}\label{eq:fix-quotient-pl}
 \Fix(\overline T)=\pi(\Fix(T));
\end{equation}
\item the inverse image of a compact interval invariant under all descended
maps is a compact interval invariant under all original maps;
\item if $T=P^rQ^s$ with $r,s>0$ and the collapsed intervals are exactly all
the nondegenerate connected components of $\Fix(T)$, then $P$ and $Q$
respect the quotient, and every common fixed point of the descended pair
lifts to a common fixed point of $P,Q$;
\item after fixed sets have become finite, a respected quotient followed by
restriction to a common invariant core cannot increase their cardinalities;
it strictly decreases the cardinality of $\Fix(T)$ whenever one quotient
fibre contains two distinct $T$-fixed points.
\item Suppose $P,Q\in\mathscr F$ have no common fixed point and admit no
proper nonempty compact interval which is forward invariant under both maps.
If $\pi$ is nonconstant, then the descended pair
$\overline P,\overline Q$ also has no common fixed point.
\end{enumerate}
\end{proposition}

\begin{proof}
The existence, continuity, and piecewise-linear regularity of the descended maps are
Lemma~\ref{lem:pl-finiteness}(5).  Uniqueness follows from the
surjectivity of $\pi$, and commutation follows after composing with $\pi$.
If $F$ is onto and $y=\pi(x)$, choose $z$ with $F(z)=x$; then
$\overline F(\pi(z))=y$, so $\overline F$ is onto.

One inclusion in \eqref{eq:fix-quotient-pl} is immediate.  Conversely,
if $\overline T(\pi(x))=\pi(x)$, the fibre
$C=\pi^{-1}(\pi(x))$ is a compact interval and $T(C)\subseteq C$.
The interval fixed-point theorem supplies a point of $C\cap\Fix(T)$.
This proves (2).  If $J'\subseteq I'$ is invariant under every descended
map, then $\pi^{-1}(J')$ is a compact interval and the intertwining identity
$\pi F=\overline F\pi$ proves (3).

For (4), put $X=\Fix(T)$.  On $X$ the restrictions of the factors have the
explicit two-sided inverses
\[
 (P|_X)^{-1}=P^{r-1}Q^s|_X,\qquad
 (Q|_X)^{-1}=Q^{s-1}P^r|_X.
\]
For example, $P^{r-1}Q^s$ preserves $X$ because it commutes with $T$, and
both composites of its restriction with $P|_X$ equal $T|_X=\id_X$; the
argument for $Q$ is identical.  Thus $P$ and $Q$ are homeomorphisms of $X$
and permute all its connected components.  In particular, each
nondegenerate component is sent onto another nondegenerate component, so
collapsing \emph{all} such components satisfies
\eqref{eq:respect-fibres}.

Let $y$ be fixed by both descended maps and let $C=\pi^{-1}(y)$.  If $C$ is
a singleton, the two intertwining identities show that its point is fixed
by both maps.  Otherwise $C$ is a nondegenerate component of $X$.
The same identities give $P(C),Q(C)\subseteq C$, while component permutation
gives equality.  On $C$ the restrictions are commuting interval
homeomorphisms.  If both preserve orientation, each fixes both endpoints.
If one reverses orientation, it has a unique fixed point in $C$;
commutation sends that point to another fixed point of the reversing map,
so the other map fixes it as well.

Finally, \eqref{eq:fix-quotient-pl} says that a quotient only identifies
fixed points, while restriction to a core only discards them.  If a fibre
contains two distinct fixed points, their images coalesce, giving the strict
inequality in (5).

For (6), suppose that $y\in I'$ is fixed by both descended maps and put
$C=\pi^{-1}(y)$.  The intertwining identities give
$P(C)\subseteq C$ and $Q(C)\subseteq C$.  If $C$ is a singleton, its point is
fixed by both original maps.  If $C$ is nondegenerate, then it is a proper
compact interval because $\pi$ is nonconstant, and it is forward invariant
under both maps.  Either alternative contradicts the hypotheses.
\end{proof}

All quotient operations below are canonical quotients of the preceding
form.  Proposition~\ref{prop:quotient-package} will be used without further
comment to retain the piecewise-linear category, preserve a counterexample after a
nonconstant respected quotient of a minimal core, compare fixed-set
cardinalities, and lift a common fixed point from a fixed-component quotient.

\subsection{Normalization and the selected repelling side}

Put
\begin{equation}\label{eq:pl-W}
 W=(PQ)^2.
\end{equation}

\begin{lemma}[Normalization and lexicographic descent]
\label{lem:pl-normalization}
Starting from a piecewise-linear counterexample on a minimal common
invariant core, one can obtain another such counterexample for which
\(X=\Fix(W)\) is finite and, on putting
\[
 \nu=\#X
\]
and
\begin{equation}\label{eq:pl-words}
 L=2(2\nu)!,\qquad A=P^L,\qquad B=Q^L,\qquad
 R=AB=W^{L/2},\qquad S=R^2=W^L,
\end{equation}
the sets \(\Fix(R)\) and \(\Fix(S)\) are finite as well.  Normalized piecewise-linear
counterexamples therefore exist, and among all of them one can choose a pair
minimizing
\begin{equation}\label{eq:pl-mintriple}
 \bigl(\#\Fix(W),\#\Fix(R),\#\Fix(S)\bigr)
\end{equation}
lexicographically.

For such a minimizing pair, a later nonconstant respected finite interval
quotient, followed by restriction to a minimal common invariant core and
fixed-component normalization, cannot lower the first fixed count.  If that
count remains \(\nu\), its canonical exponent is still \(L\), and the
descended words \(R,S\) are literally those in \eqref{eq:pl-words}; their
fixed counts cannot increase.  If the second count is also unchanged, the
third count cannot decrease.
\end{lemma}

\begin{proof}
Lemma~\ref{lem:pl-finiteness} says that \(\Fix(W)\) has finitely many
components.  If \(\Fix(W)=\I\), collapsing its sole nondegenerate component
gives the one-point quotient, whose unique point is fixed by both descended
maps; Proposition~\ref{prop:quotient-package}(4) would then lift a common
fixed point of \(P,Q\).  Hence this case is impossible, and the quotient
which collapses all nondegenerate components of \(\Fix(W)\) is nonconstant.
Collapse them and return to a minimal common invariant core.  The quotient
remains finite-piece piecewise linear, a common fixed point would lift, and
\eqref{eq:fix-quotient-pl} says that the new \(\Fix(W)\) is finite.
Every later respected quotient or core
restriction can only identify or discard points of this finite set: by
\eqref{eq:fix-quotient-pl} it cannot increase its cardinality or
recreate a nondegenerate fixed component.

We now make the three minimizations in dependency order.  Among all
minimal-core counterexamples for which \(\Fix(W)\) is finite, choose one
for which
\[
 \nu=\#\Fix(W)
\]
is least.  This class is nonempty by the preceding paragraph.  For every
pair in this first minimizing class the integer
\(L=2(2\nu)!\), and hence the words \(R=P^LQ^L\) and \(S=R^2\), have the
same prescribed form.

For a pair attaining the first minimum, \(\Fix(R)\neq\I\).  Indeed, if
\(R\) fixed the whole interval, collapsing its sole nondegenerate fixed
component would give the one-point quotient, and
Proposition~\ref{prop:quotient-package}(4) would lift the quotient's common
fixed point.  Collapse all nondegenerate components of \(\Fix(R)\), then
restrict to a minimal common invariant core.  Exact fixed-point descent
shows that the resulting \(W\)-fixed count is at most \(\nu\); by the choice
of \(\nu\) it is exactly \(\nu\).  Consequently the canonical exponent is
still \(L\), the descended \(R\) is exactly \(\overline P^L\overline Q^L\),
and its fixed set is finite.  Thus there are first-stage minimizers with
finite \(\Fix(R)\).  Among them choose one for which
\[
 \rho=\#\Fix(R)
\]
is least.

For a pair attaining these first two minima, \(\Fix(S)\neq\I\), by the
same fixed-component lifting argument.  Collapse the nondegenerate
components of \(\Fix(S)\) and return to a minimal common invariant core.
The new \(W\)-fixed count is at most \(\nu\), so it is \(\nu\); the exponent
is therefore still \(L\).  The new \(R\)-fixed count is at most \(\rho\),
so it is \(\rho\).  Exact descent now makes \(\Fix(S)\) finite.  Among the
systems attaining the first two minima and having finite \(\Fix(S)\), choose
one for which \(\#\Fix(S)\) is least.  This is precisely a lexicographic
minimum of the triple \eqref{eq:pl-mintriple}, and proves existence of the
normalization asserted in the statement.

Finally consider any later respected quotient and minimal-core restriction.
Exact fixed-point descent makes the new first count no larger than \(\nu\).
If it were smaller, fixed-component normalization of the descended pair
would produce a member of the first minimizing class with a smaller first
count, a contradiction.  Hence that count equals \(\nu\), so the canonical
exponent remains \(L\) and \(R,S\) descend literally.  Their fixed counts
cannot increase.  If the \(R\)-count fell below \(\rho\), normalizing the
descended system would contradict the second minimum.  Thus it too remains
\(\rho\); any strict decrease of the \(S\)-count then contradicts the third
minimum.  This is the lexicographic descent principle used below.
\end{proof}

Fix henceforth a normalized counterexample minimizing
\eqref{eq:pl-mintriple}.  Commutation gives \(P(X),Q(X)\subseteq X\), and on
\(X\) the factors have the explicit inverses
\begin{equation}\label{eq:pl-X-inverses}
 (P|_X)^{-1}=Q(PQ)|_X,\qquad
 (Q|_X)^{-1}=P(PQ)|_X.
\end{equation}
Indeed, both displayed composites with the indicated factor equal
\((PQ)^2=W\), which is the identity on \(X\).  Thus \(P,Q\) permute \(X\).
The set \(X\) has at least two points: if it were a singleton, both
permutations would fix its unique point.

It is useful to record the complementary words
\[
 C_P=Q(PQ),\qquad C_Q=P(PQ).
\]
Since all the factors commute, they satisfy the two-sided identities
\begin{equation}\label{eq:pl-complementary-identities}
 C_P\circ P=P\circ C_P=W,\qquad
 C_Q\circ Q=Q\circ C_Q=W,\qquad
 P\circ W=W\circ P,\qquad Q\circ W=W\circ Q.
\end{equation}

Choose consecutive points \(p<c\) of \(X\).  The continuous function
\(W-\id\) has a strict constant sign on \((p,c)\).  Reflecting the ambient
coordinate if necessary, assume
\begin{equation}\label{eq:pl-gap-sign}
 W(x)>x\qquad(p<x<c).
\end{equation}
Conjugacy by \(x\mapsto1-x\) preserves commutation, the absence of a common
fixed point, finite-piece linearity, all three fixed-set cardinalities, and
the canonical words.  Hence reflection preserves the chosen normalization
and entails no loss of generality.

A \emph{directed side} is a pair $(a,\epsilon)$ with $a\in X$ and
$\epsilon\in\{-,+\}$ for which the corresponding side of $a$ in $\I$ is
nonempty.  Write $a_+(t)=a+t$ and $a_-(t)=a-t$.  The state
$(a,\epsilon)$ is \emph{$W$-repelling} if there is $\delta>0$ such that
$W(a)=a$ and, whenever $0<t<\delta$, the point
$W(a_\epsilon(t))$ lies on that same side and
\[
 |W(a_\epsilon(t))-a|>t;
\]
here $\delta$ is chosen so that all $a_\epsilon(t)$ lie in $\I$.  By
continuity we may shrink $\delta$ so that the entire closed one-sided
interval maps into the same side.  Equation~\eqref{eq:pl-gap-sign} shows
that $(p,+)$ is repelling.
There are at most $2\nu$ repelling directed states.

\begin{lemma}[Directed-state return]\label{lem:pl-states}
The maps \(P,Q\) permute the finite set of \(W\)-repelling directed states;
on each such state their one-sided germs are strictly affine
homeomorphisms onto the image state.  Consequently \(A,B\) fix the selected
state \((p,+)\), and their selected germs are orientation-preserving \(C^1\)
local diffeomorphisms.  Their derivatives have bounded variation on a
sufficiently small selected-side interval.  They also fix \(p\) and \(c\).
\end{lemma}

\begin{proof}
Let $(a,\epsilon)$ be a repelling state.  Its one-sided neighborhood can be
shrunk so that \(W\), \(P\), and the branch of \(C_P\)
met after applying \(P\) are affine on the relevant one-sided intervals.
The germ of \(P\) cannot be constant: otherwise the factorization
\(C_P\circ P=W\) would make \(W\) constant on a smaller source-side
interval, contrary to strict motion away from the base.  Thus this affine
germ is strictly monotone and maps the source side homeomorphically onto a
one-sided interval based at \(P(a)\).  Since
\[
 W(P(a))=P(W(a))=P(a),
\]
the new base belongs to \(X\); the image side therefore defines a directed
state.

Use directed distance from the base on both sides.  In these coordinates a
strictly monotone map from one chosen side onto its image side is increasing,
whether or not it preserves the ambient orientation.  If \(t\) is close
enough to the source base, then \(W(t)\) lies farther into the same side than
\(t\).  Hence \(P(W(t))\) lies farther into the image side than \(P(t)\).
The commutation identity
\[
 W(P(t))=P(W(t))
\]
shows that the image state is again \(W\)-repelling.

The same argument, using the complementary factor $P$ in
\eqref{eq:pl-complementary-identities}, shows that $C_P$ also sends
repelling states to repelling states.  On directed states the composite
induced by \(C_P\circ P=W\) is the identity,
because \(W\) fixes the base and preserves its chosen side.  Thus the action
of \(P\) has a left inverse and is injective on the finite state set; it is
therefore a permutation.  Replacing \(P,C_P\) by
\(Q,C_Q\) proves the same assertion for \(Q\).

Every permutation of at most \(2\nu\) objects has order dividing
\((2\nu)!\).  Hence the multiple \(L=2(2\nu)!\) makes both \(P^L=A\)
and \(Q^L=B\) fix every directed state.  Their germs at the selected state
return the right side of \(p\) to itself and therefore have positive affine
slope.  They are local \(C^1\) diffeomorphisms there, and the logarithm of
each derivative is constant, hence of bounded variation.  The same exponent
fixes every point of \(X\), since \(P,Q\) permute a set of size \(\nu\).
In particular \(A,B\) fix \(p,c\).
\end{proof}

\subsection{Finite arrival and local eventual surjectivity}

Put \(J=[p,c]\).  Since \(A,B\) fix its endpoints, continuity gives
\begin{equation}\label{eq:pl-J-expands}
 J\subseteq A(J),\qquad J\subseteq B(J).
\end{equation}
Let
\[
 K_{i,j}=A^iB^j(J),\qquad
 K=\cl\bigcup_{n\geq0}K_{n,n}.
\]
Because $J\subseteq A(J)$ and $J\subseteq B(J)$ and the maps commute, the
family $K_{i,j}$ is increasing in each index.  Continuity and compactness
give
\[
 A(K)\subseteq\cl\bigcup_{n\geq0}K_{n+1,n}\subseteq K.
\]
Conversely $K_{n,n}\subseteq K_{n+1,n}=A(K_{n,n})\subseteq A(K)$ for every
$n$, and $A(K)$ is compact, so $K\subseteq A(K)$.  The same two inclusions
with $B$ in place of $A$ prove
\begin{equation}\label{eq:pl-K-invariant}
 A(K)=K=B(K).
\end{equation}

We shall also use the action on $Y=\Fix(S)$.  Commutation makes $Y$
invariant under $P,Q$ and under the following complementary words, and on
$Y$ there are explicit two-sided inverses
\begin{equation}\label{eq:pl-S-inverses}
 (P|_Y)^{-1}=P^{2L-1}Q^{2L}|_Y,\qquad
 (Q|_Y)^{-1}=Q^{2L-1}P^{2L}|_Y.
\end{equation}
Indeed, each composite in the first identity, in either order, is
$P^{2L}Q^{2L}=S$, hence the identity on $Y$; the second identity is the
same calculation.  Thus $P$ and $Q$ permute $Y$ and preserve distinctness
there.

We record the phase-intersection argument which controls every interval
quotient below.  If a compact interval \(C\) satisfies
\(P^L(C)=C=Q^L(C)\), put
\[
 D_{i,j}=P^iQ^j(C),\qquad
 (i,j)\in(\mathbb Z/L\mathbb Z)^2,
\]
and join two labels when their intervals meet.  Translation by either
coordinate is a graph automorphism.  For example, applying \(P\) proves
that an intersection implies an intersection after translating the first
coordinate.  Conversely, if
\(y\in P(D_\alpha)\cap P(D_\beta)\), then
\[
 P^{L-1}(y)\in P^L(D_\alpha)\cap P^L(D_\beta)
 =D_\alpha\cap D_\beta,
\]
because \(P^L\) fixes every phase interval.  The argument for \(Q\) is
identical.  The two translations act transitively on the label set, so this
finite interval-intersection graph is vertex transitive.  Notice that no
injectivity of \(P\) or \(Q\) on the whole interval has been used.

Apply the construction to \(C=K\).  First suppose that its phase graph is
connected.  A connected vertex-transitive interval graph is complete.  To
see this, choose a representing interval with least right endpoint.  All of
its neighbors contain that endpoint and hence form a clique, so this vertex
is simplicial.  Vertex transitivity makes every vertex simplicial.  If two
vertices were nonadjacent, a shortest path between them would have length at
least two, and the first three vertices on that path would contradict
simpliciality of the middle one.  Thus every two phase intervals meet.  The
one-dimensional Helly property now gives
\[
 C_0=\bigcap_{0\leq i,j<L}P^iQ^j(K)\neq\varnothing.
\]
Translation of the labels and the elementary inclusion
\(P(\bigcap D_{i,j})\subseteq\bigcap P(D_{i,j})\) give
\[
 P(C_0)\subseteq C_0,\qquad Q(C_0)\subseteq C_0.
\]
The intersection is a nonempty compact interval; minimality of the common
invariant core forces \(C_0=\I\).  Since \(C_0\subseteq K\), it follows that
\(K=\I\).

It remains to exclude a disconnected phase graph.  Set
\[
 F=\bigcup_{0\leq i,j<L}D_{i,j}.
\]
Graph components correspond exactly to the connected components of this
finite union of compact intervals.  The label translations show that
\(P(F)=F=Q(F)\) and permute these components.  More explicitly, a continuous
image of one component lies in a component, while the inverse label
translation proved above makes the induced maps on the finite component set
bijective.  Hence collapsing every component of \(F\) defines a nonconstant
finite interval quotient respected by both \(P\) and \(Q\).

Every phase interval contains
\[
 P^iQ^j(p)\quad\hbox{and}\quad P^iQ^j(c).
\]
These points are distinct and fixed by \(S\): the points \(p,c\) lie in
\(\Fix(W)\subseteq\Fix(S)\), and \(P,Q\) act as permutations on
\(\Fix(S)\), by \eqref{eq:pl-S-inverses}.  Thus at least one collapsed
component contains two
distinct \(S\)-fixed points.  Exact fixed-point descent strictly lowers the
\(S\)-fixed count, unless the quotient or subsequent minimal-core
restriction has already lowered the \(W\)- or \(R\)-fixed count.  In the
first case normalization contradicts the first minimum in
Lemma~\ref{lem:pl-normalization}; if the first count is unchanged, \(L\) is
unchanged, and a drop of the second count contradicts the second minimum; if
both earlier counts are unchanged, the strict drop of the third count
contradicts the final minimum.  Proposition~\ref{prop:quotient-package}(6)
ensures throughout that the descended maps remain a counterexample.  The
disconnected case is therefore impossible, and the connected case has
already shown that \(K=\I\).

Consequently
\begin{equation}\label{eq:pl-dense-diagonal}
 R^n(J)\quad(n\geq0)
\end{equation}
is an increasing sequence of intervals with dense union in \(\I\).

Write \(R^n(J)=[\ell_n,r_n]\).  Then
\(\ell_n\downarrow0\) and \(r_n\uparrow1\).  We next prove that both ambient
endpoints occur at a finite stage.  Notice that $R$ and $S$ are onto, since
$P$ and $Q$ are onto.

Suppose, for definiteness, that \(0\) never occurs.  We claim that
\begin{equation}\label{eq:pl-singleton-endpoint}
 S(0)=0,\qquad S^{-1}(0)=\{0\}.
\end{equation}
Every interior point belongs to some even diagonal interval
\(R^{2n}(J)\).  If such a point mapped to zero under \(S=R^2\), then zero
would belong to \(R^{2n+2}(J)\), contrary to the assumption.  Hence
$S^{-1}(0)\subseteq\{0,1\}$.  If the other ambient endpoint has already
occurred, it belongs to some even diagonal interval as well, so the same
argument excludes it as a preimage of zero.  Onto-ness of $S$ then forces
$S^{-1}(0)=\{0\}$, and in particular $S(0)=0$.

It remains to justify the case in which both ambient endpoints are
asymptotic.  Every \(x\in(0,1)\) lies in some diagonal interval.  If
\(R(x)\) were an ambient endpoint, that endpoint would occur at the next
stage.  Hence \(R((0,1))\subseteq(0,1)\).  Onto-ness of $R$ forces each
ambient endpoint to have a preimage in $\{0,1\}$.  There are only two
candidates, so $R|_{\{0,1\}}$ is a permutation.  Thus \(S=R^2\) fixes both
endpoints and neither has any other preimage.  This proves
\eqref{eq:pl-singleton-endpoint} in all cases.

After shrinking \(\varepsilon>0\), the map \(S\) is affine on
\([0,\varepsilon]\).  It cannot be constant there, since that would
contradict \eqref{eq:pl-singleton-endpoint}.  Its slope cannot be negative,
because \(S(0)=0\) and \(S\) takes values in \([0,1]\).  Hence the slope is
positive and \(S\) is strictly increasing on \([0,\varepsilon]\).
Compact separation gives
\[
 \eta=\min S([\varepsilon,1])>0.
\]
For all sufficiently large \(n\), both
\(\ell_{2n}<\varepsilon\) and \(\ell_{2n+2}<\eta\).  Choose
\(x_n\in R^{2n}(J)=[\ell_{2n},r_{2n}]\) with
\(S(x_n)=\ell_{2n+2}\).  The definition of \(\eta\) forces
\(x_n<\varepsilon\).  Since \(S\) is increasing on
\([0,\varepsilon]\) and \(\ell_{2n}\leq x_n\), we have
\(S(\ell_{2n})\leq\ell_{2n+2}\).  The reverse inequality follows from
\(S(\ell_{2n})\in S(R^{2n}(J))=R^{2n+2}(J)\).  Hence
\[
 \ell_{2n+2}=S(\ell_{2n}).
\]
The positive sequence on the left decreases to zero.  Since \(\Fix(S)\) is
finite, the strict sign of \(S(x)-x\) on a sufficiently small punctured
endpoint interval is therefore negative.  Thus every sufficiently small
positive point has iterates decreasing to zero.  More explicitly, after
shrinking \(\rho>0\), singleton-fibre positivity and the strict sign give
\[
 0<S(x)<x\qquad(0<x\leq\rho).
\]
Every orbit beginning in \((0,\rho]\) decreases to a fixed point in
\([0,\rho]\), necessarily zero.  If an arbitrary basin point first enters
\((0,\rho)\) after \(n\) iterates, continuity of \(S^n\) gives a neighborhood
which also enters \((0,\rho)\).  Hence the basin of zero is open.

Let \(U_0\) be the basin component containing zero and put
\(C_e=\cl(U_0)\).  This is a nondegenerate endpoint interval because the
basin contains $[0,\rho]$.  Under the standing assumption that zero never
occurs in a diagonal interval, already $0\notin J$, so $p>0$.  The point
$p$ is fixed by $S$ and is not in the basin of zero.  It is therefore an
interior barrier which the basin component beginning at zero cannot pass;
in particular $C_e$ is proper.

Write $C_e=[0,u]$.  The basin is forward invariant; the connected set
$S(U_0)$ contains zero and lies in the basin, so $S(U_0)\subseteq U_0$.
Continuity gives $S(C_e)\subseteq C_e$.  Since $u<1$ and the basin is open,
$u$ is not in the basin and $U_0=[0,u)$.  If $S(u)\in U_0$---including
$S(u)=0$---then $u$ itself would converge to zero, a contradiction.
Therefore $S(u)=u$.  The interval $S(C_e)$ contains both $0$ and $u$, and
hence
\begin{equation}\label{eq:pl-basin-onto}
 S(C_e)=C_e.
\end{equation}

We also need the elementary singleton-fibre factorization.  If commuting
onto interval maps \(U,V\) satisfy
\[
 (UV)^{-1}(0)=\{0\},\qquad UV(0)=0,
\]
then, using both factorizations \(UV,VU\),
\[
\begin{alignedat}{2}
 U^{-1}(0)&=\{V(0)\},\qquad&V^{-1}(0)&=\{U(0)\},\\
 U^{-1}(U(0))&=\{0\},&V^{-1}(V(0))&=\{0\}.
\end{alignedat}
\]
For example, if $U(x)=0$, choose $y$ with $V(y)=x$; then
$UV(y)=0$, so $y=0$ and $x=V(0)$.  Conversely $U(V(0))=UV(0)=0$.
Also, $U(x)=U(0)$ implies
$UV(x)=VU(x)=VU(0)=UV(0)=0$, whence $x=0$.  The other two identities are
symmetric.

For an onto interval map, a value having an ambient endpoint as its unique
preimage must itself be an ambient endpoint.  Indeed, if the value were in
$(0,1)$, preimages of $0$ and $1$ would both differ from the endpoint, and
the intermediate value theorem on the interval joining those preimages
would give a second preimage of the intermediate value.  Applied to the last
two singleton-fibre identities, this shows that $U(0),V(0)\in\{0,1\}$.
Moreover $U(V(0))=0$ and $V(U(0))=0$.  If $U(0)=0$, then $U$ fixes zero.
If $U(0)=1$, then $V(0)\neq0$, so $V(0)=1$ and $U(1)=0$; thus $U$
interchanges the ambient endpoints.  The same dichotomy holds for $V$.

Apply this to the two factorizations
\[
 S=P\bigl(P^{2L-1}Q^{2L}\bigr)
  =Q\bigl(P^{2L}Q^{2L-1}\bigr),
\]
whose factors are commuting and onto.  Equation~\eqref{eq:pl-singleton-endpoint}
is exactly the singleton-fibre hypothesis.  It follows that each of \(P,Q\)
either fixes zero or interchanges the ambient endpoints.  Since \(L\) is even,
\[
 A(0)=0=B(0).
\]
Commutation with $S$ sends the basin to itself, and $A(0)=B(0)=0$; hence
$A(U_0),B(U_0)\subseteq U_0$ and therefore
\(A(C_e),B(C_e)\subseteq C_e\).  Since
\(S=A^2B^2\), we have
$C_e=S(C_e)\subseteq A(C_e)\subseteq C_e$ and similarly for $B$.
Consequently
\[
 A(C_e)=C_e=B(C_e).
\]

The phase images \(P^iQ^j(C_e)\), \(0\leq i,j<L\), now obey the same
vertex-transitive intersection-graph dichotomy.  If their graph is
connected, interval Helly gives a common invariant intersection contained in
the proper interval \(C_e\).  Minimality would make this intersection equal
to $\I$, an immediate contradiction.  If the graph is disconnected, its
union cannot equal $\I$; collapse the connected components of that union to
obtain a nonconstant quotient.  Proposition~\ref{prop:quotient-package}(6) says
that the descended pair is still a counterexample.  Every phase interval
contains the two distinct
images of the endpoints of \(C_e\), which lie in \(\Fix(S)\); distinctness
is preserved by the inverses in \eqref{eq:pl-S-inverses}.
Thus the quotient strictly lowers \(\#\Fix(S)\), unless an earlier fixed
count has already fallen.  If the \(W\)-count falls, normalize again and
contradict the first minimum; if it is unchanged, the exponent \(L\) is
unchanged, and a fall of the \(R\)-count contradicts the second minimum; if
both are unchanged, the strict fall of the \(S\)-count contradicts the third
minimum.  The reflected argument treats the endpoint \(1\).  Hence no
ambient endpoint is merely asymptotic.  Choose arrival times \(n_0,n_1\)
for zero and one and put \(N=\max\{n_0,n_1\}\).  The diagonal intervals are
nested, so \(R^N(J)\) contains both ambient endpoints and therefore
\begin{equation}\label{eq:pl-finite-arrival}
 R^N(J)=\I
\end{equation}
for some \(N\).

There remains a terminal-gap issue.  Suppose \(c=1\) and \(p>0\).
For every \(x\in[p,1]\), equation~\eqref{eq:pl-gap-sign} gives
\(W(x)\geq x\), with equality only at the endpoints.  Inductively every
iterate \(W^m(x)\) remains in \([p,1]\).  Since
\(R^N=W^{NL/2}\), this gives
\[
 R^N(J)\subseteq[p,1],
\]
contradicting \eqref{eq:pl-finite-arrival}.  Thus \(c=1\) forces \(p=0\).

If \(p=0,c=1\), then consecutiveness of \(p,c\) in \(X\) gives
\(X=\{0,1\}\).  By \eqref{eq:pl-X-inverses}, the restrictions \(P|_X\) and
\(Q|_X\) are permutations; each therefore either fixes both endpoints or
interchanges them.  Suppose a factor, say \(P\), interchanges the endpoints.
Lemma~\ref{lem:pl-states} makes its selected germ a local
diffeomorphism from the right of zero to the left of one, hence locally
decreasing.  For small \(x>0\),
\[
 W(x)>x\quad\Longrightarrow\quad P(W(x))<P(x).
\]
Commutation gives \(W(P(x))<P(x)\), contradicting
\eqref{eq:pl-gap-sign}.  Hence neither factor interchanges the
endpoints; both fix zero, already a common fixed point.  In the surviving
counterexample,
\begin{equation}\label{eq:pl-c-internal}
 c<1.
\end{equation}

Finally, fix $0<\delta<c-p$, put $x_0=p+\delta$, and write
$x_n=W^n(x_0)$.  As long as $x_n\in(p,c)$,
\eqref{eq:pl-gap-sign} gives $x_{n+1}>x_n$.  Hence either some
$x_n\geq c$, or the sequence remains in $(p,c)$ and increases to a limit
$d\leq c$.  In the latter case continuity gives $W(d)=d$; since there is no
$W$-fixed point in $(p,c)$, necessarily $d=c$.

Choose $z_0,z_1\in J$ with
\[
 R^N(z_0)=0,\qquad R^N(z_1)=1.
\]
Because $R$ fixes $p,c$ and $c<1$, the point $z_1$ lies in $(p,c)$; likewise
$z_0\in(p,c)$ unless $p=0$ and $z_0=p$.  Put
\[
 r=\max\{z_0,z_1\}<c.
\]
(In the exceptional case $z_0=p$, the same strict inequality follows from
$z_1\in(p,c)$.)  The orbit argument supplies $n$ with $x_n\geq r$.
The image $W^n([p,p+\delta])$ is an interval containing both
$W^n(p)=p$ and $x_n$, and hence contains $[p,r]$, in particular $z_0,z_1$.
Applying $R^N=W^{NL/2}$ gives an interval containing zero and one; it is
therefore all of $\I$.  Since $R$ and $W$ are powers of the same map, the
composite is the single iterate $W^{n+NL/2}$.  If
$\delta\geq c-p$, the original interval contains a smaller selected-side
interval to which the preceding argument applies.  We have proved the
fundamental local property
\begin{equation}\label{eq:pl-LEO}
 \boxed{\text{For every }0<\delta\leq1-p\text{ there is }k\text{ such that }
 W^k([p,p+\delta])=\I.}
\end{equation}

In particular, suppose $F$ and $G$ commute with $W$ and agree on a
selected-side neighborhood of $p$.  By continuity they agree on
$[p,p+\delta]$ for some $\delta>0$.  Choose $k$ from
\eqref{eq:pl-LEO}.  Given $y\in\I$, choose
$x\in[p,p+\delta]$ with $W^k(x)=y$.  Then
\[
 F(y)=F(W^k(x))=W^k(F(x))=W^k(G(x))=G(W^k(x))=G(y).
\]
Thus $F=G$ on $\I$.  We call this \emph{germ determination}.

\subsection{The one-sided \texorpdfstring{\(C^{1+\mathrm{BV}}\)}{C1+BV} orbit circle}

For an orientation-preserving \(C^1\) local diffeomorphism, write
\(C^{1+\mathrm{BV}}\) when \(\log Df\) has bounded variation on a sufficiently
small closed one-sided interval.  This class is closed under composition and
inverse, because
\[
 \log D(g\circ h)=(\log Dg)\circ h+\log Dh,
 \qquad
 \log D(g^{-1})=-(\log Dg)\circ g^{-1}.
\]
Composition with a monotone map does not increase variation.  Moreover, on
every compact representative the positive continuous derivative has range in
some interval \([c,C]\subset(0,\infty)\).  The restrictions of \(\log\) and
\(\exp\) to the corresponding compact ranges are Lipschitz, and therefore
\begin{equation}\label{eq:pl-log-derivative-BV}
 \Var(\log Df)<\infty\quad\Longleftrightarrow\quad \Var(Df)<\infty.
\end{equation}
Thus this convention is exactly strong enough for the derivative-BV
hypothesis in Denjoy's theorem below.
Every selected germ occurring here is affine after the domain is
shrunk sufficiently.  Once its positive slope is known to be nonzero, its
logarithmic derivative is constant, so it is automatically
\(C^{1+\mathrm{BV}}\).

We use three standard one-dimensional facts.  Navas's
bounded-variation Kopell lemma \cite[Theorem 4.1.1]{Navas} says that if
commuting orientation-preserving diffeomorphisms $u,v$ of $[0,1)$ are
respectively $C^{1+\mathrm{BV}}$ and $C^1$, if $u$ has no interior fixed
point, and if $v$ has one, then $v=\id$.  The Poincar\'e translation-number
criterion \cite[Propositions 2.2.2 and 2.2.3]{Navas} says that an
orientation-preserving circle-homeomorphism lift $G$ has translation number
$m/n\in\mathbb Q$, $n>0$, exactly when $G^n(t)=t+m$ for some $t$.
Finally, Denjoy's theorem
\cite[Theorem 3.1.1]{Navas} says that an orientation-preserving $C^1$ circle
diffeomorphism whose derivative has bounded variation and whose rotation
number is irrational is topologically conjugate to the corresponding rigid
rotation.

We first record the local form of Kopell's lemma that will be used below.

\begin{lemma}[Local BV Kopell corollary]\label{lem:local-BV-Kopell}
Let \(u,v\) be commuting orientation-preserving one-sided \(C^1\) local
diffeomorphism germs fixing zero.  Suppose that \(u\) is
\(C^{1+\mathrm{BV}}\), has no fixed point on the punctured side, and that
\(v\) has a punctured-side fixed point.  Then \(v\) is the identity germ.
\end{lemma}

\begin{proof}
Replace \(u\) by \(u^{-1}\), if necessary, so that \(u(x)<x\).  Choose a
fixed point \(y>0\) of \(v\) in a common representative domain on which
\(u,v\), their local inverses, and the commutation identity are all defined.
Then every
\(u^j(y)\), \(j\geq0\), is fixed by \(v\), and \(u^j(y)\downarrow0\).
Extend \(u\) above \(y\) to a fixed-point-free
\(C^{1+\mathrm{BV}}\) diffeomorphism of \([0,1)\), preserving the already
given germ.  For completeness, write \(z=u(y)<y\) and \(r=Du(y)>0\).
We construct the extension explicitly.  Put \(a=y-z>0\),
\(s=1-r<1\), and choose
\(0<\eta<\min\{r,1\}/2\).  Fix \(\kappa>0\), and choose
\(Y\in(y,1)\) so close to \(y\) that the function
\[
 q(x)=s+\frac{x-y}{Y-y}(-\kappa-s),\qquad
 d(x)=a+\int_y^x q(t)\,dt \quad (y\leq x\leq Y),
\]
is positive.  Since both endpoint values of \(q\) are at most
\(1-\eta\), so is \(q\) throughout this interval.  Write
\(a_Y=d(Y)>0\) and set
\(\lambda=\kappa(1-Y)^2/a_Y\).  On \([Y,1)\), continue \(d\) by
\[
 d(x)=a_Y\exp\!\left[-\lambda
   \left(\frac1{1-x}-\frac1{1-Y}\right)\right].
\]
The value and first derivative match at \(Y\), because the right derivative
there is \(-\lambda a_Y/(1-Y)^2=-\kappa\).  For the variation calculation,
put \(w=(1-x)^{-1}\).  On the tail,
\[
 d'=-\lambda w^2d,\qquad
 d''=d(\lambda^2w^4-2\lambda w^3).
\]
Since \(d=C e^{-\lambda w}\) there, where
\(C=a_Ye^{\lambda/(1-Y)}\), and \(dx=dw/w^2\),
\[
 \int_Y^1|d''(x)|\,dx
 \leq C\int_{(1-Y)^{-1}}^\infty
 e^{-\lambda w}(\lambda^2w^2+2\lambda w)\,dw<\infty.
\]
Thus \(d'\) has finite variation on the tail.  It equals the affine
function \(q\) on \([y,Y]\), and there is no jump at \(Y\), so \(d'\) has
bounded variation on all of \([y,1)\).  The same formula shows that the tail
and all its derivatives tend to zero at one.  Consequently \(d\) is positive
and \(C^1\), \(d'\leq1-\eta\), and \(d(x)\to0\) as \(x\uparrow1\).  Define
\(u(x)=x-d(x)\) for \(x\geq y\).
Its value and derivative at \(y\) are \(z\) and \(r\), so it glues to the
given representative.  Moreover \(Du=1-d'\geq\eta\), \(u(x)<x\), and
\(u(x)\to1\).  Thus the result is an orientation-preserving diffeomorphism
of \([0,1)\) onto itself with no interior fixed point.  Since \(Du\) has
bounded variation and is bounded away from zero, \(\log Du\) has bounded
variation as well.  This is the required extension.

Keep \(v\) on \([0,y]\).  On the successive fundamental intervals above
\(y\), extend it by
\[
 v\bigm|_{u^{-k}([u(y),y])}
   =u^{-k}\bigl(v\bigm|_{[u(y),y]}\bigr)u^k,
 \qquad k=1,2,\ldots .
\]
Here \(a_k=u^{-k}(y)\) increases to one, and the \(k\)-th displayed formula
acts on \([a_{k-1},a_k]\).  It fixes both endpoints, so the pieces fit
continuously and define an increasing homeomorphism of \([0,1)\).  Indeed,
the increasing sequence \((a_k)\) can have no limit below one, since such a
limit would be an interior fixed point of \(u\); hence these fundamental
intervals exhaust \([y,1)\).
At the common endpoint \(a_k\), the chain rule gives
\[
 D(u^{-k}vu^k)(a_k)=Dv(y),\qquad
 D(u^{-(k+1)}vu^{k+1})(a_k)=Dv(u(y)),
\]
because the derivatives of each inverse pair cancel at the corresponding
orbit point.  Since
\(v(y)=y\), differentiating \(uv=vu\) at \(y\) gives
\[
 Du(y)Dv(y)=Dv(u(y))Du(y),
\]
and \(Du(y)>0\), so these two one-sided derivatives agree.  Thus the pieces
glue \(C^1\).  The conjugation formulas also give \(uv=vu\) on each
fundamental interval, and continuity gives the identity at every seam.
Only finitely many seams meet a compact subinterval of \([0,1)\), so the
extension is a \(C^1\) diffeomorphism there and still fixes \(y\).
The BV Kopell lemma makes it the identity, so the original \(v\) is the
identity germ.
\end{proof}

\begin{lemma}[Orbit-circle dichotomy]\label{lem:orbit-circle}
Let \(f,A,B\) be pairwise commuting orientation-preserving
\(C^{1+\mathrm{BV}}\) one-sided local diffeomorphism germs fixing zero, and
suppose
\[
 f(x)>x\quad(0<x\ll1),\qquad AB=f^d,\qquad d\in\mathbb Z_{>0}.
\]
The physical germ \(A\) selects a real lift translation number \(\alpha\),
normalized by \(\tau_f(f)=1\).  Exactly one of the following holds.
\begin{enumerate}
\item If \(\alpha=m/n\in\mathbb Q\), where \(n>0\), then
\begin{equation}\label{eq:rational-germ-powers}
 A^n=f^m,\qquad B^n=f^{dn-m}
\end{equation}
as germs.
\item If \(\alpha\notin\mathbb Q\), there are an increasing homeomorphism
\(\chi:(0,\varepsilon)\to(-\infty,T)\) and simultaneous germ identities,
valid on a sufficiently far-left tail,
\begin{equation}\label{eq:irrational-germ-translations}
 \chi f\chi^{-1}(s)=s+1,\quad
 \chi A\chi^{-1}(s)=s+\alpha,\quad
 \chi B\chi^{-1}(s)=s+d-\alpha.
\end{equation}
\end{enumerate}
\end{lemma}

\begin{proof}
\begin{claim}[Direct-limit orbit line]\label{clm:direct-limit-orbit-line}
For the germs \(f,A,B\) above there is a Hausdorff ordered line
\(\mathscr L\), carrying a \(C^{1+\mathrm{BV}}\) atlas, in which the physical
punctured side is a neighborhood of the negative end.  The three germs
extend uniquely by equivariance to pairwise commuting
\(C^{1+\mathrm{BV}}\) orientation-preserving diffeomorphisms of
\(\mathscr L\), and all their germ relations remain valid.  The extension
of \(f\) acts freely, \(\mathscr L/\langle f\rangle\) is a
\(C^{1+\mathrm{BV}}\) circle, and there is a
\(C^{1+\mathrm{BV}}\) coordinate \(H:\mathscr L\to\mathbb R\) satisfying
\(HfH^{-1}(s)=s+1\).
\end{claim}

\begin{proof}[Proof of the claim]
Shrink the representatives so that $f,A,B$ and their inverses are defined on
the intervals used below.  Here is an explicit model for the direct limit.
For a sufficiently small representative interval \(U=(0,\varepsilon)\),
take classes \([x,j]\), \(x\in U\), \(j\in\mathbb Z\), under the
equivalence relation generated by
\begin{equation}\label{eq:direct-limit-relation}
 [f(x),j]=[x,j+1]
\end{equation}
whenever the displayed terms are defined.  The ordered direct limit is an
open line \(\mathscr L\).  The physical punctured side embeds by
\(x\mapsto[x,0]\).  After shrinking \(U\), the selected local inverse
satisfies \(0<f^{-1}(x)<x\) throughout \(U\).  Hence all its forward iterates
are defined, and \(f^{-k}(x)\) decreases to a limit \(\ell\geq0\).
Continuity gives \(f(\ell)=\ell\); the strict germ inequality excludes a
positive such limit, so \(\ell=0\).  Thus the embedded physical side is a
neighborhood of
the negative end.  Using a sufficiently small representative of each class,
the germ extends to the free deck map

\[
 f_{\mathscr L}[x,j]=[f(x),j]=[x,j+1].
\]

We continue to denote this extension by \(f\).  Choose \(x_0>0\) and put
\(x_1=f(x_0)>x_0\).  Every point of \(\mathscr L\) has a representative
\([x,j]\) with \(x\in D=[x_0,x_1]\); the only boundary identifications are
\([x_1,j]=[x_0,j+1]\).  Thus \(\mathscr L\) is the ordered gluing of the
copies \(D_j=D\times\{j\}\).  Reduction to \(D\) is unique except for these
two seam representatives, and
\[
 H_0([x,j])=j+\frac{x-x_0}{x_1-x_0}\qquad(x\in D)
\]
is well defined at each seam.  Its restriction to every copy is a
homeomorphism onto \([j,j+1]\), and the piecewise inverse has the quotient
topology.  Hence \(H_0\) is a homeomorphism
\(\mathscr L\cong\mathbb R\); in particular the quotient is Hausdorff and
the asserted ordered open line is genuine.

For regularity, enlarge each copy slightly across its two endpoints using
\eqref{eq:direct-limit-relation}.  The overlap maps are restrictions of
iterates of \(f\) and \(f^{-1}\), so they form a
\(C^{1+\mathrm{BV}}\) atlas compatible with this topology.  The deck map
is represented in these charts by \(f\) and its conjugates.  Its quotient
\[
 \mathscr O_f=\mathscr L/\langle f\rangle
\]
is a topological circle.  A fundamental interval is
$D=[x_0,x_1]$, with $x_0$ identified with $x_1$ by $f$.

We spell out the regularity at this seam.  Choose a positive $C^1$ function
$\rho$ on $D$, with derivative of bounded variation, satisfying
\[
 \rho(x_1)Df(x_0)=\rho(x_0).
\]
For example, prescribe the two positive endpoint values in this ratio and
join them by a positive affine function.  Rescale \(\rho\) so that
\(\int_{x_0}^{x_1}\rho=1\), and
set
\[
 h(x)=\int_{x_0}^x\rho(t)\,dt,\qquad x\in D.
\]
Extend $h$ to a coordinate $H:\mathscr L\to\mathbb R$ by

\[
 H([x,j])=h(x)+j \qquad (x\in D).
\]

The values agree at the seam because \(h(x_1)=1\) and \(h(x_0)=0\).
To check the derivative, use the \(j\)-th physical coordinate on both sides.
Just to the right of the seam, write the point as
\([x,j+1]=[f(x),j]\).  The derivative there, in the coordinate \(f(x)\),
tends to \(\rho(x_0)/Df(x_0)\); from the left it tends to \(\rho(x_1)\).
These are equal by the endpoint condition on \(\rho\).  Hence there is no
derivative jump.  Each side is \(C^{1+\mathrm{BV}}\), and a compact part of
\(\mathscr L\) meets only finitely many copies, so the total variation on
that compact part is finite.  Positivity of \(\rho\) gives the same
statement for the local inverse.  Thus \(H\) is a
\(C^{1+\mathrm{BV}}\) coordinate and
\begin{equation}\label{eq:orbit-deck}
 HfH^{-1}(s)=s+1.
\end{equation}
It follows in particular that $\mathscr O_f$ carries the standard
$C^{1+\mathrm{BV}}$ circle structure.

For \(G=A\) or \(G=B\), define its equivariant extension by
\begin{equation}\label{eq:direct-limit-G}
 G_{\mathscr L}[x,j]=[G(x),j],
\end{equation}
using a sufficiently small representative of the class.  This is
independent of the representative: relation
\eqref{eq:direct-limit-relation} and \(Gf=fG\) give

\[
 [G(f(x)),j]=[f(G(x)),j]=[G(x),j+1].
\]

The same construction with the inverse germ proves that
\(G_{\mathscr L}\) is an orientation-preserving diffeomorphism, and the
definition gives \(G_{\mathscr L}f=fG_{\mathscr L}\).  In overlapping
direct-limit charts its formulas are conjugates \(f^kGf^{-k}\).  Adjacent
formulas have equal values by commutation, and differentiating \(Gf=fG\)
shows that their first derivatives agree at the seam.  All such formulas
are \(C^{1+\mathrm{BV}}\).  Since a compact interval in \(\mathscr L\)
meets only finitely many charts, its derivative variation is a finite sum
of the piecewise variations, with no seam jump.  Thus \(G_{\mathscr L}\)
is \(C^{1+\mathrm{BV}}\) on the line.
Formula \eqref{eq:direct-limit-G} also preserves every germ relation:
evaluate both words on a representative sufficiently close to zero.  In
particular the extensions of \(A,B,f\) commute pairwise and satisfy
\(AB=f^d\).  Equivariance and agreement on the embedded negative end make
the extensions unique.  This proves the claim.
\end{proof}

Choose the line, extensions, and coordinate supplied by
Claim~\ref{clm:direct-limit-orbit-line}.  For \(G\in\{A,B\}\), put
\(\widetilde G=HG_{\mathscr L}H^{-1}\).  Commutation with \(f\) gives the
exact deck identity
\[
 \widetilde G(s+1)=\widetilde G(s)+1.
\]
Consequently \(D\widetilde G(s+1)=D\widetilde G(s)\); restricting to
\([0,1]\) and identifying its endpoints produces an orientation-preserving
\(C^{1+\mathrm{BV}}\) circle diffeomorphism.  This proves the required
regularity of the circle maps induced by \(A\) and \(B\).

The actual action of $G$ on $\mathscr L$, rather than only its quotient
circle map, selects this lift $\widetilde G$ commuting with the deck
translation $T_1(s)=s+1$.  A different lift $T_1^k\widetilde G$ would
represent the different germ $f^kG$.  Define
\[
 \alpha=\tau_f(A)=\lim_{r\to\infty}
       \frac{\widetilde A^{\,r}(s)-s}{r}.
\]
This is independent of $s$, its class modulo $\mathbb Z$ is the rotation
number of the induced circle map, and the germ identity $AB=f^d$ gives the
exact lift identity
\begin{equation}\label{eq:selected-lift-product}
 \widetilde A\widetilde B=T_1^d.
\end{equation}

Suppose first that $\alpha=m/n\in\mathbb Q$, with $n>0$.  The
Poincar\'e theorem provides $s$ such that
$\widetilde A^n(s)=s+m$.  Translating $s$ by a sufficiently negative integer
places the corresponding point in the physical representative of the germ;
the equality is unchanged because the lifts commute with $T_1$.  Hence
$H_0=A^nf^{-m}$ has a punctured-side fixed point.  It commutes with the
fixed-point-free $C^{1+\mathrm{BV}}$ germ $f$, so
Lemma~\ref{lem:local-BV-Kopell} gives $H_0=\id$.  Therefore
$A^n=f^m$.  Since $B=f^dA^{-1}$, also
$B^n=f^{dn-m}$.

Now suppose that $\alpha$ is irrational.  Denjoy's theorem conjugates the
circle map induced by $A$ to rotation by $\alpha\bmod\mathbb Z$.  Lift a
degree-one conjugacy to an increasing homeomorphism of universal covers
which commutes with $T_1$.  The conjugated selected lift of $A$ is a
translation.  Its displacement is exactly $\alpha$, rather than merely its
residue class, because translation number is invariant under degree-one lift
conjugacy.  Equation~\eqref{eq:selected-lift-product} then sends the selected
lift of $B$ to translation by $d-\alpha$, while
\eqref{eq:orbit-deck} sends $f$ to translation by one.  Restricting this
lifted coordinate to a sufficiently far-left part of the physical end gives
an increasing homeomorphism
$\chi:(0,\varepsilon)\to(-\infty,T)$ and the three identities in
\eqref{eq:irrational-germ-translations}.
\end{proof}

Apply Lemma~\ref{lem:orbit-circle} in the physical selected-side coordinate
\(x-p\) to \(f=W\) and to \(A,B\), where \(d=L/2\) and \(AB=W^d\).
There is also an elementary one-map dilation chart
\[
 \jmath:[0,\varepsilon_0)\longrightarrow[p,p+\delta_0),\qquad
 W(\jmath(t))=\jmath(et)\quad(0\leq t<\varepsilon_0/e).
\]
Here is the construction in full.  Choose \(x_0>p\) so close to \(p\) that
\(W\) is an increasing local diffeomorphism on \([p,x_0]\), that
\(x_1=W(x_0)>x_0\), and that the selected inverse branch
\(g:[p,x_1]\to[p,x_0]\) is defined.  Choose an increasing homeomorphism
\(\eta:[1,e]\to[x_0,x_1]\) with \(\eta(1)=x_0\) and \(\eta(e)=x_1\).
For \(0<t<e\), choose \(n\geq0\) such that
\[
 e^{-n}\leq t\leq e^{\,1-n}
\]
and put
\[
 \jmath(t)=g^n\bigl(\eta(e^nt)\bigr);
\]
also put \(\jmath(0)=p\).  At the common endpoint of two consecutive
parameter intervals the two formulas agree, since
\[
 g^{n+1}(\eta(e))=g^{n+1}(x_1)
 =g^n(x_0)=g^n(\eta(1)).
\]
The intervals \(g^n([x_0,x_1])\) are consecutive and their endpoints tend
to \(p\): the decreasing sequence \(g^n(x_0)\) has a fixed-point limit, and
the selected germ has no fixed point other than \(p\).  Hence the displayed
pieces fit to an increasing homeomorphism
\(\jmath:[0,e)\to[p,x_1)\), continuous also at zero.  Finally, for
\(0<t<1\), choosing \(n\geq1\) with
\(e^{-n}\leq t\leq e^{1-n}\) gives
\[
 W(\jmath(t))=g^{n-1}(\eta(e^nt))=\jmath(et),
\]
and the identity at \(t=0\) follows from \(W(p)=p\).  Thus one may take
\(\varepsilon_0=e\) and \(\delta_0=x_1-p\).

Suppose first that \(\alpha=m/n\) is rational.  If \(m<0\), write
\(h=-m>0\).  Choose \(0<\varepsilon<\varepsilon_0/e\) so small that the
selected inverse germ of \(W\), all the identities used below, and the chart
formula
\[
 W^{-h}(\jmath(t))=\jmath(e^{-h}t)
\]
are defined on \(0\leq t\leq\varepsilon\).  Put
\(I_\varepsilon=\jmath([0,\varepsilon])\).  Here \(W^{-h}\) means only the
corresponding iterate of this selected local inverse; no global inverse of
\(W\) is being used.  Since \(W^{-h}(I_\varepsilon)\subseteq I_\varepsilon\),
induction using \(A^n=W^{-h}\) on this chart gives, for every \(k\geq0\),
\begin{equation}\label{eq:pl-rational-shrinking}
 A^{nk}(I_\varepsilon)
 =W^{-hk}(I_\varepsilon)
 =\jmath([0,e^{-hk}\varepsilon]).
\end{equation}
Thus these compact intervals shrink to \(\{p\}\), and every use of the germ
identity occurs inside its chosen domain.  Choose \(q\) from
\eqref{eq:pl-LEO} with \(W^q(I_\varepsilon)=\I\).  Surjectivity of \(A\) and
commutation give, for every \(k\),
\[
 \I=A^{nk}(\I)
    =A^{nk}(W^q(I_\varepsilon))
    =W^q(A^{nk}(I_\varepsilon)).
\]
On the other hand, continuity of the fixed map \(W^q\), together with
\(W^q(p)=p\), gives
\[
 \sup_{x\in A^{nk}(I_\varepsilon)}|W^q(x)-p|\longrightarrow0.
\]
The last displayed set therefore cannot equal the nondegenerate interval
\(\I\) for every \(k\), a contradiction.  If \(m=0\), germ determination
gives \(A^n=\id\)
globally, so \(P^{Ln}=\id\).  Then \(P\) is a finite-order interval
homeomorphism: if it preserves orientation it is the identity, while if it
reverses orientation its unique fixed point is preserved by \(Q\).  Either
case gives a common fixed point.  Applying the same argument to
\(B^n=W^{dn-m}\) shows that \(dn-m\) is neither negative nor zero.  Hence
\(0<m<dn\), and
\[
 A^{\,n(dn-m)}=B^{\,nm}.
\]
Germ determination makes this equality global, that is,
\(P^{Ln(dn-m)}=Q^{Lnm}\).  Both powers are positive, so
Lemma~\ref{lem:pl-common-iterate} below again gives a common fixed
point of \(P,Q\).  Thus the rational alternative is impossible in the
standing counterexample.

We are in the irrational alternative.  Exponentiate its translation
coordinate and return to the physical side.  This gives the germ chart
\[
 \iota:[0,e^T)\longrightarrow\I,\qquad
 \iota(0)=p,\qquad
 \iota(t)=p+\chi^{-1}(\log t)\quad(0<t<e^T),
\]
and set
\[
 \lambda=e,\qquad a=e^\alpha,\qquad b=e^{d-\alpha},\qquad
 \Lambda=e^d=\lambda^d.
\]
Then, as selected-side germ identities (that is, whenever \(t\) and the
scaled arguments displayed below lie in \([0,e^T)\)),
\begin{equation}\label{eq:actor-dilations}
 W(\iota(t))=\iota(\lambda t),\qquad
 A(\iota(t))=\iota(at),\qquad
 B(\iota(t))=\iota(bt),
\end{equation}
and
\begin{equation}\label{eq:pl-scale-product}
 ab=\Lambda>1.
\end{equation}

No onto centralizer return can have scale below one.  If, for example,
\(0<a<1\), choose \(\varepsilon>0\) so small that the \(A\)-germ identity is
valid throughout \(I_\varepsilon=\iota([0,\varepsilon])\).  Because its
image is nested inside the same chart, induction gives
\begin{equation}\label{eq:pl-subunit-shrinking}
 A^k(I_\varepsilon)=\iota([0,a^k\varepsilon])
 \qquad(k\geq0).
\end{equation}
Choose \(q\) with \(W^q(I_\varepsilon)=\I\).  Surjectivity and commutation
then give
\[
 \I=A^k(\I)=A^k(W^q(I_\varepsilon))
    =W^q(A^k(I_\varepsilon))
 \qquad(k\geq0).
\]
The intervals in \eqref{eq:pl-subunit-shrinking} shrink to \(\{p\}\), so
continuity of the fixed map \(W^q\) gives
\[
 \sup_{x\in A^k(I_\varepsilon)}|W^q(x)-p|\longrightarrow0.
\]
Again the right side cannot equal \(\I\) for every \(k\), a contradiction.
The identical argument excludes \(0<b<1\).
Here \(a=1\) would make \(\alpha=0\), and \(b=1\) would make \(\alpha=d\);
both contradict irrationality.  Thus
\begin{equation}\label{eq:pl-scales-positive}
 a,b>1.
\end{equation}

We shall use the Fugate--McLean fixed-set theorem
\cite[Corollary 1.5, p.~613]{FugateMcLean}, whose exact conclusion is that
the fixed-point set of a periodic homeomorphism of a tree-like continuum is
a nonvoid continuum.

\begin{lemma}[Common positive iterate]
\label{lem:pl-common-iterate}
Let \(F,G\colon\I\to\I\) be commuting continuous maps.  If
\[
 F^m=G^n
\]
for some positive integers \(m,n\), then \(F,G\) have a common fixed point.
\end{lemma}

\begin{proof}
Put \(T=F^m=G^n\) and
\[
 J_\infty=\bigcap_{r\geq0}T^r(\I).
\]
The sets in this intersection are nested nonempty compact intervals.  Also
\begin{equation}\label{eq:pl-eventual-range-onto}
 T(J_\infty)=J_\infty.
\end{equation}
Indeed, one inclusion is immediate.  For \(y\in J_\infty\), the nonempty
compact sets
\[
 T^{-1}(y)\cap T^r(\I),\qquad r\geq0,
\]
are nested: nonemptiness follows from \(y\in T^{r+1}(\I)\).  Their
intersection contains a preimage of \(y\) in \(J_\infty\), proving the other
inclusion.  Commutation gives
\[
 F(J_\infty)\subseteq J_\infty,\qquad
 G(J_\infty)\subseteq J_\infty.
\]
If \(J_\infty\) is a singleton, it is already common fixed.  Assume it is
nondegenerate.

Form the natural extension
\[
 X_T=\{(x_0,x_1,\ldots)\in J_\infty^{\mathbb N_0}:
       T(x_{i+1})=x_i\text{ for every }i\geq0\}.
\]
By \eqref{eq:pl-eventual-range-onto}, this is a nonempty compact
connected metric space.  For completeness, in the product
\(J_\infty^{\mathbb N_0}\) let
\[
 X_T^{(N)}=\{x:T(x_{i+1})=x_i\text{ for }0\leq i<N\}.
\]
Once \(x_N\) and the unconstrained coordinates after \(N\) are chosen, the
coordinates before \(N\) are determined.  Thus \(X_T^{(N)}\) is homeomorphic
to a product of compact intervals, and in particular is nonempty, compact,
and connected.  These sets are nested, and
\(X_T=\bigcap_NX_T^{(N)}\); the nested-intersection theorem therefore proves
the assertion.  Give \(X_T\) the metric
\[
 d_X(x,y)=\sum_{i=0}^{\infty}2^{-i-1}|x_i-y_i|.
\]
The coordinate projection \(\pi_N\colon X_T\to J_\infty\) is onto: prescribe
the \(N\)-th coordinate, determine the earlier ones by forward iteration,
and choose the later ones recursively using
\eqref{eq:pl-eventual-range-onto}.  If
\(\pi_N(x)=\pi_N(y)\), the bonding equations make \(x_i=y_i\) for
\(0\leq i\leq N\), and hence
\[
 d_X(x,y)\leq\sum_{i>N}2^{-i-1}=2^{-N-1}.
\]
Thus, for every \(\varepsilon>0\), some \(\pi_N\) is an
\(\varepsilon\)-map onto an interval, meaning that every point fibre has
diameter less than \(\varepsilon\).  This is precisely the defining
approximation property for an arc-like continuum.  Since an arc is a finite
tree, every arc-like continuum is tree-like.

The shift
\[
 \Sigma(x_0,x_1,\ldots)=(T(x_0),x_0,x_1,\ldots)
\]
is a homeomorphism, with inverse deleting the zeroth coordinate.
Coordinatewise application defines commuting maps \(\widehat F,\widehat G\)
of \(X_T\), and
\[
 \widehat F^m=\Sigma=\widehat G^n.
\]
A continuous selfmap whose positive power is a homeomorphism is itself a
homeomorphism: injectivity follows from injectivity of the power, surjectivity
from surjectivity of the power, and a continuous bijection of a compact space
onto a Hausdorff space is a homeomorphism.  Thus \(\widehat F,\widehat G\)
are commuting homeomorphisms.

Write
\[
 m=d_0m_0,\qquad n=d_0n_0,\qquad \gcd(m_0,n_0)=1,
\]
and choose integers \(r_0,s_0\) with
\[
 r_0n_0+s_0m_0=1.
\]
Negative powers are now legitimate.  Set
\[
 U=\widehat F^{r_0}\widehat G^{s_0},\qquad
 V=\widehat F^{m_0}\widehat G^{-n_0}.
\]
Then \(V^{d_0}=\id\), and the B\'ezout identity gives
\begin{equation}\label{eq:pl-bezout-split}
 \widehat F=U^{n_0}V^{s_0},\qquad
 \widehat G=U^{m_0}V^{-r_0}.
\end{equation}

The quoted Fugate--McLean theorem makes
\[
 \mathcal C=\Fix(V)
\]
a nonempty compact connected subset of \(X_T\).  Since \(U\) commutes with
\(V\), it maps \(\mathcal C\) into itself.  We prove directly that
\(U|_{\mathcal C}\) has a fixed point.  Given \(\varepsilon>0\), choose
\(N\) with \(2^{-N-1}<\varepsilon\).  The set
\(\pi_N(\mathcal C)=[\alpha,\beta]\) is an interval.  If
\(\alpha=\beta\), every \(z\in\mathcal C\) satisfies
\(d_X(U(z),z)<\varepsilon\).  Otherwise choose
\(z_\alpha,z_\beta\in\mathcal C\) over the two endpoints.  The continuous
function
\[
 z\longmapsto\pi_N(U(z))-\pi_N(z)
\]
is nonnegative at \(z_\alpha\) and nonpositive at \(z_\beta\).  Its image of
the connected set \(\mathcal C\) therefore contains zero.  In either case
there is \(z_\varepsilon\in\mathcal C\) with
\[
 d_X(U(z_\varepsilon),z_\varepsilon)<\varepsilon.
\]
Compactness, followed by \(\varepsilon\downarrow0\), supplies
\(z\in\mathcal C\) with \(U(z)=z\).  Since also \(V(z)=z\),
\eqref{eq:pl-bezout-split} gives
\(\widehat F(z)=z=\widehat G(z)\).  The zeroth coordinate is a common fixed
point of \(F,G\).
\end{proof}

If \(a^m=b^n\) for positive integers \(m,n\), the dilation coordinate makes
\(A^m\) and \(B^n\) equal as germs.  Germ determination makes them equal
globally:
\[
 P^{Lm}=Q^{Ln}.
\]
 Lemma~\ref{lem:pl-common-iterate} then gives a common fixed point of
 \(P,Q\).
Consequently the sole remaining case is
\begin{equation}\label{eq:pl-independent-scales}
 a,b>1,\qquad \frac{\log a}{\log b}\notin\mathbb Q.
\end{equation}

\subsection{The canonical ray and its tail core}

For \(t\geq0\), choose \(n\) large enough that \(t/\lambda^n\) lies in the
domain of \(\iota\), and define
\begin{equation}\label{eq:pl-ray}
 \Phi(t)=W^n\left(\iota\left(\frac{t}{\lambda^n}\right)\right).
\end{equation}
Whenever \(n\) is large enough that both germ arguments are admissible,
\[
 W^{n+1}\left(\iota\left(\frac{t}{\lambda^{n+1}}\right)\right)
 =W^n\left(\iota\left(\frac{t}{\lambda^n}\right)\right).
\]
Indeed, this is one application of the \(W\)-identity in
\eqref{eq:actor-dilations}, followed by \(W^n\).  Induction makes
\eqref{eq:pl-ray} independent of every sufficiently large choice of \(n\).
One \(n\) works on each compact
parameter interval, so \(\Phi\) is continuous; near zero it equals
\(\iota\).  Property~\eqref{eq:pl-LEO} makes \(\Phi\) onto.  Commutation
and the germ identities give
\begin{equation}\label{eq:pl-ray-equivariance}
 \Phi(at)=A(\Phi(t)),\qquad
 \Phi(bt)=B(\Phi(t)),\qquad
 \Phi(\Lambda t)=R(\Phi(t)).
\end{equation}

For \(T>0\), put
\[
 \mathcal C_T=\cl\bigl(\Phi([T,\infty))\bigr),\qquad
 \Omega=\bigcap_{T>0}\mathcal C_T.
\]
The set \(\Phi([T,\infty))\) is nonempty and connected, so its closure in
the compact interval \(\I\) is a nonempty compact interval.  Thus each
\(\mathcal C_T\) is a compact interval.  The family decreases with \(T\), so
\(\Omega\) is itself a nonempty compact interval.  Equivariance gives
\(A(\mathcal C_T)=\mathcal C_{aT}\).  Here is the closure bookkeeping.
If \(x=\lim_k\Phi(t_k)\in\mathcal C_T\), with \(t_k\geq T\), continuity and
equivariance give
\(A(x)=\lim_k\Phi(at_k)\in\mathcal C_{aT}\).  Conversely, if
\(y=\lim_k\Phi(s_k)\in\mathcal C_{aT}\), take \(s_k\geq aT\) and put
\(t_k=s_k/a\).  A subsequence of \(\Phi(t_k)\) converges, by compactness of
\(\I\), to some \(x\in\mathcal C_T\), and then
\(A(x)=\lim_k\Phi(s_k)=y\).  This proves the asserted equality.  Hence
\(A(\Omega)\subseteq\Omega\).  Conversely, if \(y\in\Omega\), then the
nested compact sets
\[
 \mathcal C_T\cap A^{-1}(y)\qquad(T>0)
\]
are nonempty because \(y\in\mathcal C_{aT}=A(\mathcal C_T)\).  Their
intersection contains \(x\in\Omega\) with \(A(x)=y\).  Thus
\(A(\Omega)=\Omega\), and the identical argument for \(B\) gives
\begin{equation}\label{eq:pl-Omega-invariant}
 A(\Omega)=\Omega=B(\Omega).
\end{equation}

Suppose first that \(\Omega=\{q\}\).  If \(x\) is fixed by \(R\), choose
\(t\) with \(\Phi(t)=x\).  If \(t>0\), then
\[
 x=R^n(x)=\Phi(\Lambda^nt)\longrightarrow q;
\]
the lift \(t=0\) gives \(x=p\).  Hence
\[
 \Fix(R)\subseteq\{p,q\}.
\]
The adjacent point \(c\) is distinct from \(p\) and is fixed by \(A,B\),
hence by \(R=AB\).  Since \(\Phi(0)=p\), every lift of \(c\) is positive;
the preceding argument therefore gives \(c=q\).
The sign of \(R-\id\) is positive on \((p,q)\) and negative on
\((q,1]\): it is positive near \(p\), and there is no further fixed point
on either interval.

Commutation gives \(P(p),Q(p)\in\Fix(R)\), so each is either \(p\) or \(q\).
If \(P(p)=q\), Lemma~\ref{lem:pl-states} says that its selected germ is
nonzero and maps the right side of \(p\) to one of the two sides at \(q\).
On the source, \(R\) moves points away from \(p\).
Commutation would therefore make the target side repelling.  But on the left
of \(q\), \(R(x)>x\) moves points toward \(q\), while on the right,
\(R(x)<x\) does the same.  This order contradiction proves \(P(p)=p\).
Similarly \(Q(p)=p\), contrary to the assumption of no common fixed point.

Suppose instead that \(\Omega\) is nondegenerate and proper.  Its phase
images
\[
 P^iQ^j(\Omega),\qquad0\leq i,j<L,
\]
are nondegenerate: if one were a singleton, applying
\(P^{L-i}Q^{L-j}\) would make
\(P^LQ^L(\Omega)=R(\Omega)=\Omega\) a singleton.  They have a
vertex-transitive intersection graph by
\eqref{eq:pl-Omega-invariant} and the phase-intersection argument
above.  We now use exactly the phase-graph construction following
Lemma~\ref{lem:pl-normalization}, the quotient package in
Proposition~\ref{prop:quotient-package}, and the fixed-set identity
\eqref{eq:fix-quotient-pl}; all quotient classes below are the connected
components prescribed by that construction.  If the
graph is connected, interval Helly gives a proper common invariant
intersection.  If it is disconnected, collapse the connected components of
the union.  For every phase interval \(D_{i,j}=P^iQ^j(\Omega)\), commutation
and \(R(\Omega)=\Omega\) give
\[
 R(D_{i,j})=P^iQ^j(R(\Omega))=D_{i,j}.
\]
It follows that \(R\) maps each connected component of their union onto
itself.  Every continuous onto selfmap \(f\) of a
nondegenerate interval has at least two fixed points of \(f^2\).  Indeed, if
\(f\) has at least two fixed points there is nothing to prove.  Otherwise
let \(z\) be its unique fixed point.  It is interior: if, for example, it
were the left endpoint, then \(f(x)<x\) at every other point, whereas a
preimage of the right endpoint, supplied by surjectivity, would satisfy
\(f(x)>x\); the right-endpoint case is symmetric.  Moreover,
\(f(x)>x\) to its left while \(f(x)<x\) to its right.  If the interval is
\([a_0,b_0]\), surjectivity gives \(v<z<u\) with
\(f(v)=b_0\) and \(f(u)=a_0\).  Choose \(x_0\in[v,z)\) with \(f(x_0)=u\).
Then
\[
 f^2(x_0)-x_0=a_0-x_0\leq0,\qquad
 f^2(a_0)-a_0\geq0.
\]
Equality, or the intermediate value theorem, gives an \(f^2\)-fixed point
to the left of \(z\), in addition to \(z\).  Here \(f=R|_C\) on a component
\(C\), so \(f^2=S|_C\) by \eqref{eq:pl-words}.  Thus every collapsed
component contains at least two points of \(\Fix(S)\).  Formula
\eqref{eq:fix-quotient-pl} identifies their images with the fixed points of
the descended \(S\); consequently the quotient strictly lowers the third
count in \eqref{eq:pl-mintriple}, unless restriction and renormalization have
already lowered an earlier count.  Lemma~\ref{lem:pl-normalization} keeps the
word \(S\) unchanged whenever the first count is unchanged, while
Proposition~\ref{prop:quotient-package}(6) preserves the counterexample.
Either alternative contradicts lexicographic minimality.

We conclude that
\begin{equation}\label{eq:pl-Omega-full}
 \Omega=\I.
\end{equation}
Equivalently, every late ray tail has dense image.

\subsection{Dense tails imply exactness}

Recall that an interval map \(T\) is \emph{topologically exact} if every
nonempty relatively open set \(U\subseteq\I\) satisfies
\(T^n(U)=\I\) for some \(n\geq0\).

We first produce a positive boundary lift of the base value \(p\).
If \(p>0\), fix \(T>0\).  Density of \(\Phi([T,\infty))\), and the fact that
\(0<p<c\leq1\), give parameters \(t_-,t_+>T\) for which
\(\Phi(t_-)<p<\Phi(t_+)\).  On the compact segment with endpoints
\(t_-\) and \(t_+\), a boundary point \(t_*>T\) of the relatively open set
\(\{t:\Phi(t)>p\}\) satisfies
\begin{equation}\label{eq:pl-boundary-lift}
 \Phi(t_*)=p
\end{equation}
and, by its definition as a boundary of that open set, is approached by
values on the selected right side.

Suppose \(p=0\).  If \(\Phi^{-1}(0)=\{0\}\), then
\(R^{-1}(0)=\{0\}\): for \(R(x)=0\), choose \(x=\Phi(t)\); then
\(\Phi(\Lambda t)=0\), forcing \(t=0\).  For each parameter \(\zeta\) in
the chart domain, put
\[
 I_\zeta=\iota([0,\zeta]).
\]
Choose \(\eta>0\) so that \(\iota\) is defined on
\([0,\Lambda\eta]\) and
\[
 R(\iota(t))=\iota(\Lambda t)\qquad(0\leq t\leq\eta).
\]
Since \(R^{-1}(0)=\{0\}\), compactness gives
\[
 \delta:=\min_{x\in[\iota(\eta),1]}R(x)>0.
\]
Choose \(\varepsilon>0\) so small that
\(\varepsilon<\Lambda\eta\) and \(\iota(\varepsilon)<\delta\).
On \([0,\iota(\eta)]\) the displayed dilation identifies the full preimage
of \(I_\zeta\), while the compact-complement bound excludes every other
preimage.  Consequently, for \(0<\zeta\leq\varepsilon\),
\begin{equation}\label{eq:pl-local-inverse}
  R^{-1}(I_\zeta)=I_{\zeta/\Lambda}.
\end{equation}
If \(\Phi(t)\in I_\varepsilon\), then
\[
 \Phi(t)=R^N\bigl(\Phi(t/\Lambda^N)\bigr).
\]
Iterating \eqref{eq:pl-local-inverse} gives
\[
 \Phi(t/\Lambda^N)\in I_{\varepsilon/\Lambda^N}.
\]
For large \(N\), the ray equals the initial chart at \(t/\Lambda^N\), so
strict monotonicity of \(\iota\) yields \(t\leq\varepsilon\).  Conversely,
\(\Phi=\iota\) on a sufficiently small initial interval.  After reducing
\(\varepsilon\), therefore,
\[
 \Phi^{-1}(I_\varepsilon)=[0,\varepsilon],
\]
contradicting density of every late tail.  Thus the zero fibre has a positive
point \(t_0\).  It cannot contain an entire tail, because every late tail is
dense in the nondegenerate interval \(\I\).  Choose \(s>t_0\) with
\(\Phi(s)>0\).  A boundary in \([t_0,s]\) between the closed zero fibre and
its open complement gives \(t_*>0\) with \(\Phi(t_*)=0\), approached by
positive, and hence selected-side, values.

Equivariance and the positive selected germs show that every
\begin{equation}\label{eq:pl-actor-boundaries}
 a^rb^s t_*,\qquad r,s\in\mathbb Z_{\geq0},
\end{equation}
has the same boundary property.

We now prove exactness of \(R\).  Let \(U\) be nonempty and relatively open,
and fix any \(T>0\).  Tail density supplies \(t_0>T\) with
\(\Phi(t_0)\in U\).  Since \(\Phi^{-1}(U)\) is relatively open, it contains a
nondegenerate compact parameter interval
\[
 J_0=[u,v],\qquad T<u<v,\qquad\Phi(J_0)\subseteq U.
\]
By \eqref{eq:pl-independent-scales}, the multiplicative group
\[
 \{a^rb^s:r,s\in\mathbb Z\}
\]
is dense in \((0,\infty)\): after taking logarithms it becomes the additive
subgroup generated by \(\log a\) and \(\log b\), which is dense because their
ratio is irrational.  Choose integers \(r,s\) such that
\[
 z=t_*a^rb^s\in(u,v).
\]
For \(k\geq\max\{0,-r,-s\}\), the point
\[
 \Lambda^kz=t_*a^{r+k}b^{s+k}
\]
is a selected-side boundary lift in the interior of \(\Lambda^kJ_0\).
Consequently
\[
 R^k(\Phi(J_0))=\Phi(\Lambda^kJ_0)
\]
contains a nondegenerate selected-side interval based at \(p\).  A further
power of \(R=W^d\) maps that interval onto \(\I\): round the exponent in
\eqref{eq:pl-LEO} upward to a multiple of \(d\), using surjectivity of
\(W\).  Hence \(R\) is topologically exact.

We prove exactness of \(A\) separately.  Given a nonempty relatively open
\(U\) and any threshold \(T>t_*\), apply the preceding open-preimage argument
to obtain a compact interval as late as desired:
\[
 J=[u,v],\qquad T<u<v,\qquad\Phi(J)\subseteq U.
\]
Let \(j\) be least with \(z=\Lambda^jt_*>u\).  Then
\[
 u<z\leq\Lambda u<\Lambda v.
\]
The set of \(\rho\in(1,\Lambda)\) for which
\[
 \rho u<z<\rho v
\]
contains a nonempty compact interval \(K_0\) in its interior.  For every
\(\rho\in K_0\), the interval \(\rho J\) contains the boundary lift \(z\)
in its interior.  Thus \(\Phi(\rho J)\) contains a nondegenerate selected
interval based at \(p\).  To make the uniformity explicit, define
\[
 m(\rho)=\max\{\Phi(t):t\in[\rho u,\rho v]\}.
\]
The change of variables \(t=\rho((1-s)u+sv)\) puts this maximum over the
fixed compact set \(s\in[0,1]\).  Uniform continuity of \(\Phi\) on the
compact union of these intervals therefore makes \(m\) continuous on
\(K_0\).  The boundary property at \(z\), together with
\(z\in(\rho u,\rho v)\), gives \(m(\rho)>p\) for every \(\rho\in K_0\).
Consequently
\(\eta_0=\min_{\rho\in K_0}(m(\rho)-p)>0\).  Since
\(\Phi([\rho u,\rho v])\) is an interval containing both \(p\) and
\(m(\rho)\), we have
\begin{equation}\label{eq:pl-uniform-selected}
 [p,p+\eta_0]\subseteq\Phi(\rho J)
 \qquad(\rho\in K_0).
\end{equation}
Choose \(q\) with \(R^q([p,p+\eta_0])=\I\).  Since \(R\) is onto, the same
holds for every exponent at least \(q\).

Write uniquely
\[
 a^n=\Lambda^{k_n}\rho_n,\qquad1\leq\rho_n<\Lambda.
\]
The residues \(\rho_n\) are dense in the multiplicative circle obtained by
identifying the endpoints of \([1,\Lambda]\), because
\[
 \frac{\log a}{\log\Lambda}
 =\frac{\log a}{\log a+\log b}
\]
is irrational.  Equivalently, the fractional parts of
\(n\log a/\log\Lambda\) are dense modulo one; every tail of this sequence is
dense.  Also
\(k_n=\lfloor n\log a/\log\Lambda\rfloor\to\infty\).  We may therefore
choose a single \(n\) with \(\rho_n\in K_0\) and \(k_n\geq q\).
Then
\[
\begin{split}
 A^n(\Phi(J))
 &=\Phi(a^nJ)
 =R^{k_n}(\Phi(\rho_nJ))\\
 &\supseteq R^{k_n}([p,p+\eta_0])=\I.
\end{split}
\]
Since \(\Phi(J)\subseteq U\), this proves exactness of \(A\).  Replacing
\(a\) by \(b\) proves exactness of \(B\).  Finally \(A=P^L\) and \(B=Q^L\),
so \(P,Q\) are exact as well.  In particular, none of
\(P,Q,A,B,R\) has a flat interval.

\subsection{A common Parry coordinate and the polygonal completion}

By Lemma~\ref{lem:pl-finiteness}, the flat-free finite-piece piecewise-linear maps
\(P,Q,R\) have finitely many strict monotonicity laps and are finite-to-one.
We use the following combination of Parry's conjugacy theorem
\cite{Parry} and Shultz's published scaled-measure results
\cite[Propositions 3.6 and 4.3 and Corollaries 4.4 and 4.6]{Shultz}.  A
piecewise-monotone interval map is
\emph{essentially injective} when there are no two disjoint nondegenerate
intervals with the same image.

\medskip
\noindent\textbf{Parry--Shultz scaled-measure theorem.}
\emph{Let $T$ be a topologically transitive piecewise-monotone map with
finitely many laps on a compact interval.  Then it admits a nonatomic
full-support Borel
probability measure $\mu$ and a number
$s=\exp(h_{\rm top}(T))\geq1$ such that
\[
 \mu(T(E))=s\mu(E)
\]
for every Borel set $E$ on which $T$ is injective.  Its distribution function
conjugates $T$ to a finite-piece piecewise-linear map whose nonzero slopes
have absolute value $s$.  If $T$ is not essentially injective, the scaled
probability is unique.  Moreover, \(s>1\) if and only if \(T\) is not
essentially injective.}
\medskip

This theorem applies to $R$.  Exactness implies transitivity.  An exact
interval map cannot be monotone: an onto flat-free monotone map is a
homeomorphism, and
no iterate of a sufficiently small open interval under a homeomorphism is
the whole interval.  Hence $R$ has a genuine turn.  At a genuine turn the
two adjacent strict laps have subintervals with the same nondegenerate image,
so $R$ is not essentially injective in Shultz's sense.  Let $\mu$ be the
unique scaled probability and let $s>1$ be its scale.

For a flat-free finite-piece piecewise-monotone map \(T\), whose finitely many
monotonicity laps are therefore strict, and a bounded Borel function
\(\varphi\), define the preimage transfer operator by
\[
 (\mathcal L_T\varphi)(y)=\sum_{Tx=y}\varphi(x).
\]
If \(T\) has \(M\) laps, assign every common lap endpoint to one adjacent lap,
obtaining pairwise disjoint Borel lap sets \(I_1,\ldots,I_M\).  Each
restriction \(T|_{I_j}\) is an injective Borel map between standard Borel
spaces.  The Lusin--Souslin theorem makes \(T(I_j)\) Borel and makes the
inverse branch \(b_j:T(I_j)\to I_j\) Borel.  Therefore
\[
 \mathcal L_T\varphi
 =\sum_{j=1}^M
   \mathbf1_{T(I_j)}\,(\varphi\circ b_j)
\]
is Borel, and
\(\|\mathcal L_T\varphi\|_\infty\leq M\|\varphi\|_\infty\).  For a finite
Borel measure \(\nu\), define a set function on Borel sets by
\begin{equation}\label{eq:pl-transfer-adjoint-definition}
 (\mathcal L_T^*\nu)(E)
 =\int_{\I}\sum_{Tx=y}\mathbf 1_E(x)\,d\nu(y).
\end{equation}
Monotone convergence proves countable additivity, and the preceding fibre
bound proves finiteness.  Thus this is a finite Borel measure; equivalently,
for every bounded Borel \(\varphi\),
\[
 \int_{\I}\varphi\,d(\mathcal L_T^*\nu)
 =\int_{\I}\mathcal L_T\varphi\,d\nu.
\]
The raw preimage sum need not preserve continuity at critical values, which
is why the adjoint has been defined directly on Borel measures.  Regrouping
finite fibres, now on bounded Borel functions, gives
\begin{equation}\label{eq:pl-transfer-composition}
 \mathcal L_{U\circ T}=\mathcal L_U\mathcal L_T.
\end{equation}

\begin{lemma}[Transfer to a centralizer]\label{lem:pl-transfer}
If \(T\) is an onto, flat-free, piecewise-monotone map with finitely many laps
and commuting with \(R\), then there is a constant \(c_T>0\) such that
\[
 \mathcal L_T^*\mu=c_T\mu.
\]
\end{lemma}

\begin{proof}
First partition \(\I\) into finitely many pairwise disjoint Borel lap sets
\(I_1,\ldots,I_N\), assigning every lap endpoint to one adjacent set, so
that \(R\) is injective on each \(I_j\).  Since \(\mu\) is nonatomic, the
finite endpoint assignment is immaterial.  For every Borel set \(E\), each
image \(R(E\cap I_j)\) is Borel by the same Lusin--Souslin argument, and
\[
 \mathcal L_R\mathbf1_E
 =\sum_{j=1}^N\mathbf1_{R(E\cap I_j)}.
\]
The scaled-set identity therefore gives
\[
\begin{split}
 (\mathcal L_R^*\mu)(E)
 &=\int\mathcal L_R\mathbf1_E\,d\mu\\
 &=\sum_{j=1}^N\mu(R(E\cap I_j))
 =s\sum_{j=1}^N\mu(E\cap I_j)
 =s\mu(E).
\end{split}
\]
Conversely, this eigenmeasure identity implies the scaled-set identity on
every injectivity set.  Thus scaling is equivalent to
\begin{equation}\label{eq:pl-R-eigenmeasure}
 \mathcal L_R^*\mu=s\mu.
\end{equation}

Put \(\nu_T=\mathcal L_T^*\mu\).  If \(T\) has \(M\) laps, then
\(\mathcal L_T\mathbf1\leq M\), so \(\nu_T\) is finite.  Onto-ness gives
\(\mathcal L_T\mathbf1\geq1\), so it is nonzero.  It is nonatomic because
\[
 \nu_T(\{x\})=\mu(\{T(x)\})=0.
\]
It has full support: every nonempty open interval contains a smaller interval
\(J\) on which \(T\) is strictly monotone, and
\[
 \nu_T(J)=\mu(T(J))>0.
\]

Commutation and \eqref{eq:pl-transfer-composition} yield
\[
 \mathcal L_R^*\nu_T
 =\mathcal L_T^*\mathcal L_R^*\mu
 =s\nu_T.
\]
After division by its positive total mass, \(\nu_T\) is another nonatomic
full-support scaled probability for \(R\) with factor \(s\).  Uniqueness
gives \(\nu_T=c_T\mu\).
\end{proof}

Apply Lemma~\ref{lem:pl-transfer} to \(T=P,Q\), and write
\[
 \mathcal L_P^*\mu=C\mu,\qquad
 \mathcal L_Q^*\mu=D\mu
\]
with \(C,D>0\).  Define the distribution coordinate
\begin{equation}\label{eq:pl-distribution}
 \theta(x)=\mu([0,x]).
\end{equation}
Full support and atomlessness make \(\theta\) an increasing homeomorphism of
\(\I\).

If \(J=[x,y]\) lies in a strict monotonicity lap of \(P\), the transfer
identity gives
\[
 \mu(P(J))=C\mu(J),
\]
and hence
\[
 |\theta(P(y))-\theta(P(x))|
 =C|\theta(y)-\theta(x)|.
\]
It follows that \(\theta P\theta^{-1}\) is affine with slope \(+C\) or
\(-C\) on each of finitely many laps.  Thus it is a finite-piece PL map.
The same argument makes \(\theta Q\theta^{-1}\) finite-piece PL.

These two conjugated maps commute.

Put
\[
 p_\theta=\theta(p).
\]
From now on we use the same letters \(P,Q,A,B,R\) for the maps conjugated by
\(\theta\).  Exactness and commutation are preserved by conjugacy.  Thus
\(P,Q,A,B,R\) are exact, flat-free, finite-piece piecewise-linear maps, with
\[
 A=P^L,\qquad B=Q^L,\qquad R=AB.
\]
Set
\begin{equation}\label{term:eq:transformed-scales}
 \widehat a=C^L,\qquad
 \widehat b=D^L,\qquad
 \widehat s=\widehat a\widehat b.
\end{equation}
Every nonzero slope of \(A,B,R\) has absolute value
\(\widehat a,\widehat b,\widehat s\), respectively.  Since the original
selected germs of \(A,B\) preserve the selected side, their conjugated germs
at \(p_\theta\) have positive slopes \(\widehat a,\widehat b\); consequently
the selected germ of \(R\) has positive slope \(\widehat s\).

\subsubsection{Multiplicative independence survives the Parry coordinate}

The original selected germs of \(A,B\) move every sufficiently close point
strictly away from \(p\).  Since \(\theta\) preserves order, their conjugated
affine germs move every sufficiently close point strictly away from
\(p_\theta\).  Their positive slopes therefore satisfy
\(\widehat a,\widehat b>1\), and hence \(C,D>1\).

If \(\widehat a^{\,m}=\widehat b^{\,n}\) for some \(m,n>0\), the selected
affine germs of \(A^m\) and \(B^n\) have the same base point and the same
positive slope, so they agree on a nondegenerate selected-side interval.
Both maps commute with the exact map \(R\).  If \(J\) is a nonempty open
subinterval on which they agree, choose \(N\) with \(R^N(J)=\I\).  For
\(y=R^N(x)\), \(x\in J\), commutation gives
\[
 A^m(y)=R^N(A^m(x))=R^N(B^n(x))=B^n(y).
\]
Thus exactness transports the local equality to the whole interval, and
\(A^m=B^n\), equivalently
\(P^{Lm}=Q^{Ln}\).  Lemma~\ref{lem:pl-common-iterate} then gives a
common fixed point of \(P,Q\), contrary to the standing counterexample
assumption.  Therefore
\begin{equation}\label{term:eq:new-independent}
 \widehat a,\widehat b>1,\qquad
 \frac{\log\widehat a}{\log\widehat b}\notin\mathbb Q.
\end{equation}

\subsection{The unit-speed ray and a common Markov partition}

\subsubsection{Rebuilding the ray}

Using the positive right germ of the constant-slope map \(R\), define, for
\(t\geq0\),
\begin{equation}\label{term:eq:Psi}
 \Psi(t)=R^n\left(p_\theta+\frac{t}{\widehat s^{\,n}}\right)
\end{equation}
for any sufficiently large \(n\).  The definition stabilizes because
\[
\begin{split}
 R^{n+1}\left(p_\theta+\frac{t}{\widehat s^{\,n+1}}\right)
 &=R^n\left(
     R\left(p_\theta+\frac{t}{\widehat s^{\,n+1}}\right)\right)\\
 &=R^n\left(p_\theta+\frac{t}{\widehat s^{\,n}}\right),
\end{split}
\]
where \(R(p_\theta+x)=p_\theta+\widehat s x\) on the selected germ.
Induction proves independence of all sufficiently large \(n\), and one such
\(n\) works on every compact parameter interval.  Thus \(\Psi\) is continuous
and locally finite-piece PL.  Exactness of \(R\) excludes a constant interval
for every \(R^n\), since all later images of such an interval would remain
singletons.  On each affine piece the absolute slope of \(R^n\) is therefore
\(\widehat s^{\,n}\), so the corresponding affine piece of \(\Psi\) has
absolute slope
\(\widehat s^{\,n}/\widehat s^{\,n}=1\).

The ray is onto.  Indeed, choose a small selected-side interval contained in
the germ of \(R\).  Exactness sends its relative interior onto the whole
interval under some iterate, and \eqref{term:eq:Psi} places that image in the
range of \(\Psi\).  Near zero,
\begin{equation}\label{term:eq:Psi-initial}
 \Psi(t)=p_\theta+t.
\end{equation}
Commutation and the selected germs give
\begin{equation}\label{term:eq:Psi-equivariance}
\begin{split}
 \Psi(\widehat a t)&=A(\Psi(t)),\\
 \Psi(\widehat b t)&=B(\Psi(t)),\\
 \Psi(\widehat s t)&=R(\Psi(t)).
\end{split}
\end{equation}

\subsubsection{Lap-vertex values}

For an onto flat-free finite-piece PL map $T$, let
\[
 E(T)=\{0,1\}\cup\operatorname{Turn}(T),\qquad
 \mathcal V(T)=T(E(T)),
\]
where $\operatorname{Turn}(T)$ is the finite set of genuine changes of
monotonicity.

\begin{lemma}[Vertex values under composition]
\label{lem:vertex-composition}
For onto flat-free finite-piece PL maps $U,T$,
\begin{equation}\label{term:eq:vertex-composition}
 \mathcal V(U\circ T)=\mathcal V(U)\cup U(\mathcal V(T)).
\end{equation}
\end{lemma}

\begin{proof}
Let $x$ be a turn of the inner map $T$.  The two adjacent branches of $T$
approach $T(x)$ from the same side.  On that side $U$ has a nonzero affine
germ, even when $T(x)$ is itself a turn coordinate of $U$.  Thus the inner
turn cannot cancel in $U\circ T$.  Domain endpoints remain domain endpoints,
so
$U(\mathcal V(T))\subseteq\mathcal V(U\circ T)$.

\begin{claim}[Transverse preimage]\label{clm:pl-transverse-preimage}
For every \(y\in(0,1)\), there is \(x\in T^{-1}(y)\) at which \(T\) is
locally one-to-one.
\end{claim}

\begin{proof}
Choose \(x_0,x_1\) with \(T(x_0)=0\) and \(T(x_1)=1\).  On the segment
joining them, \(T-y\) has opposite signs at the two ends.  Because \(T\) has
finitely many nonconstant affine pieces, \(T^{-1}(y)\) is finite.  Therefore,
as one traverses this segment, at least one zero \(x\) of \(T-y\) is a
sign-changing zero.  The two one-sided affine slopes at such a zero have the
same nonzero sign.  Hence \(T\) is strictly monotone, and thus one-to-one, on
a neighborhood of \(x\).  This proves the claim.
\end{proof}

Let \(y\) be an interior turn coordinate of \(U\).  At the preimage supplied
by Claim~\ref{clm:pl-transverse-preimage}, composition with the locally
one-to-one map \(T\) preserves the outer turn of \(U\).  Thus \(U(y)\) is a
vertex value of \(U\circ T\).  If \(y\in\{0,1\}\), every interior point of
\(T^{-1}(y)\) is a global extremum; flat-freeness makes it a genuine turn.
Hence a preimage of \(y\) is either a domain endpoint or a turn of \(T\),
and the first paragraph shows again that \(U(y)\) is a vertex value of the
composite.  This accounts for \(U(E(U))\) and proves
\(\mathcal V(U)\subseteq\mathcal V(U\circ T)\).

Conversely, if $x\notin E(T)$, then $T$ is locally strictly monotone and
one-to-one at $x$; a turn of $U\circ T$ at $x$ can therefore occur only when
$T(x)$ is a turn of $U$.  If $x\in E(T)$, its composite value lies in
$U(\mathcal V(T))$.  These alternatives prove the reverse inclusion and
hence \eqref{term:eq:vertex-composition}.
\end{proof}

Define
\begin{equation}\label{term:eq:TA-TB}
 \mathcal T_A(m)=\bigcup_{0\leq i<m}A^i(\mathcal V(A)),\qquad
 \mathcal T_B(n)=\bigcup_{0\leq j<n}B^j(\mathcal V(B)).
\end{equation}
Iterating Lemma~\ref{lem:vertex-composition} and comparing $A^mB^n$ with
$B^nA^m$ gives, for all $m,n\geq1$,
\begin{equation}\label{term:eq:rectangle-identity}
 \mathcal T_A(m)\cup A^m(\mathcal T_B(n))
 =
 \mathcal T_B(n)\cup B^n(\mathcal T_A(m)).
\end{equation}

\subsubsection{Finiteness of the postvertex sets}

\begin{lemma}[Finite postvertex sets]
\label{lem:finite-postvertex}
The sets
\[
 \mathcal T_A=\bigcup_{i\geq0}A^i(\mathcal V(A)),\qquad
 \mathcal T_B=\bigcup_{j\geq0}B^j(\mathcal V(B))
\]
are finite.
\end{lemma}

\begin{proof}
Put
$\mathcal T=\bigcup_{m\geq1}\mathcal T_A(m)$ and
$\mathcal U=\mathcal V(B)$.  Taking $n=1$ in
\eqref{term:eq:rectangle-identity} gives
\[
 \mathcal T_A(m)\cup A^m(\mathcal U)
 =\mathcal U\cup B(\mathcal T_A(m)),
\]
and hence
\begin{equation}\label{term:eq:postvertex-differences}
\begin{aligned}
 \mathcal T_A(m)\setminus B(\mathcal T_A(m))&\subseteq\mathcal U,\\
 B(\mathcal T_A(m))\setminus\mathcal T_A(m)&\subseteq A^m(\mathcal U).
\end{aligned}
\end{equation}
\begin{claim}[Finite exceptional set]\label{clm:pl-postvertex-exceptional}
There is a finite set \(Z\) of \(A\)-periodic points such that
\begin{equation}\label{term:eq:B-T}
 B(\mathcal T)\subseteq\mathcal T\cup Z.
\end{equation}
\end{claim}

\begin{proof}
Let \(y\in B(\mathcal T)\setminus\mathcal T\), and choose \(m_0\) so that a
\(B\)-preimage of \(y\) belongs to \(\mathcal T_A(m_0)\).  For every
\(m\geq m_0\), the second inclusion in
\eqref{term:eq:postvertex-differences} gives \(y\in A^m(\mathcal U)\).
Because \(\mathcal U\) is finite, there are \(u\in\mathcal U\) and
\(m_0\leq i<j\) with
\[
 A^i(u)=y=A^j(u).
\]
It follows that \(A^{j-i}(y)=y\).  For a fixed \(u\), every periodic point
on its forward orbit belongs to its unique eventual cycle, when that cycle
exists.  The union \(Z\) of these cycles over the finite set \(\mathcal U\)
is finite, consists of \(A\)-periodic points, and contains every such \(y\).
This proves the claim.
\end{proof}

Suppose, toward a contradiction, that \(\mathcal T\) is infinite.

\begin{claim}[Action on infinite tail classes]
\label{clm:pl-tail-class-action}
The points of \(\mathcal T\) with infinite forward \(A\)-orbit form a
nonempty union of finitely many tail-equivalence classes, and \(B\) induces
a selfmap of this finite class set.
\end{claim}

\begin{proof}
Declare \(x\sim y\) when their forward \(A\)-orbits meet.  Since
\(\mathcal T\) is generated by the finite seed \(\mathcal V(A)\), every
infinite orbit in \(\mathcal T\) is tail-equivalent to the orbit of one of
these seeds.  Hence there are only finitely many classes, and at least one
exists because \(\mathcal T\) is infinite.

If \(w\) has infinite \(A\)-orbit, then \(B(w)\) also has infinite
\(A\)-orbit.  Otherwise
\[
 B(A^i(w))=A^i(B(w))
\]
would send the infinitely many distinct points \(A^i(w)\) into a finite set,
contrary to finite-to-one-ness of \(B\).  Claim
\ref{clm:pl-postvertex-exceptional} then puts \(B(w)\) in \(\mathcal T\),
because every point of \(Z\) has a finite \(A\)-orbit.  Finally, if
\(A^i(x)=A^j(y)\), commutation gives
\(A^i(B(x))=A^j(B(y))\).  Thus the induced class map is well defined.
\end{proof}

\begin{claim}[Global tail relation]\label{clm:pl-tail-relation-global}
There are \(r\geq1\) and \(p_0,q_0\geq0\) such that
\begin{equation}\label{term:eq:global-tail-relation}
 A^{p_0}B^r=A^{q_0}
\end{equation}
on \(\I\).
\end{claim}

\begin{proof}
A selfmap of a nonempty finite set has a periodic element.  Choose a tail
class of period \(r\geq1\) for the class map in Claim
\ref{clm:pl-tail-class-action}, and choose \(w\) in that class.  The points
\(B^r(w)\) and \(w\) are tail-equivalent, so there are \(p_0,q_0\geq0\)
such that
\begin{equation}\label{term:eq:tail-relation}
 A^{p_0}B^r(w)=A^{q_0}(w).
\end{equation}
Commutation propagates this equality to every point \(A^i(w)\).  These
points are distinct, and the two finite-piece PL maps
\(A^{p_0}B^r\) and \(A^{q_0}\) have a finite common affine partition.
One open cell therefore contains at least two agreement points, so the two
affine restrictions agree throughout a nondegenerate open interval \(J_1\).

Both maps commute with the exact map \(R\).  Choose \(N\) with
\(R^N(J_1)=\I\).  If \(y=R^N(x)\), \(x\in J_1\), then
\[
 A^{p_0}B^r(y)=R^N A^{p_0}B^r(x)
              =R^N A^{q_0}(x)=A^{q_0}(y).
\]
This proves the global identity.
\end{proof}

\begin{claim}[Slope contradiction]\label{clm:pl-tail-slope-contradiction}
The identity in Claim~\ref{clm:pl-tail-relation-global} contradicts
\eqref{term:eq:new-independent}.
\end{claim}

\begin{proof}
At a point outside the finite breakpoint sets, equality of the two maps
implies equality of their absolute slopes:
\[
 \widehat a^{p_0}\widehat b^r=\widehat a^{q_0}.
\]
If \(p_0=q_0\), this says \(\widehat b^r=1\), contrary to
\(\widehat b>1\).  If \(p_0<q_0\), it says
\(\widehat b^r=\widehat a^{q_0-p_0}\), contrary to multiplicative
independence.  If \(p_0>q_0\), it says
\(\widehat a^{p_0-q_0}\widehat b^r=1\), although both factors exceed one.
All cases are impossible.
\end{proof}

The contradiction proves that \(\mathcal T=\mathcal T_A\) is finite.  The
same argument with \(A,B\) interchanged proves that \(\mathcal T_B\) is
finite.
\end{proof}

\subsubsection{A common finite Markov set}

\begin{lemma}[Common finite Markov set]
\label{lem:common-Markov-set}
There is a finite set $\mathcal K\subseteq\I$ containing $0,1,p_\theta$
such that
\[
 A(\mathcal K)\subseteq\mathcal K,\qquad
 B(\mathcal K)\subseteq\mathcal K,
\]
and both $A$ and $B$ are affine and monotone between consecutive points of
$\mathcal K$.
\end{lemma}

\begin{proof}
Once $\mathcal T_A(m)$ has stabilized to $\mathcal T_A$,
\eqref{term:eq:postvertex-differences} applies for every larger $m$.  For
$u\in\mathcal V(B)$ whose forward $A$-orbit is eventually periodic, let
$C_A(u)$ denote its eventual cycle, and set
\[
 Z_A=\bigcup_{\substack{u\in\mathcal V(B)\\
                 u\text{ has eventually periodic }A\text{-orbit}}}C_A(u).
\]
The argument proving \eqref{term:eq:B-T} shows that
$B(\mathcal T_A)\setminus\mathcal T_A\subseteq Z_A$.  Thus $Z_A$ is finite
and consists of $A$-periodic points.  If \(Z_A=\varnothing\), put \(r_A=1\);
otherwise choose \(r_A\) divisible by the periods of all its points.  Then
\(Z_A\subseteq\Fix(A^{r_A})\).  The latter fixed set is finite: on every
affine cell the slope of \(A^{r_A}\) has absolute value
\(\widehat a^{r_A}>1\), so \(A^{r_A}-\id\) has nonzero slope and at most one
zero.  There are only finitely many cells.

The set \(\mathcal T_A\) is forward invariant under \(A\), and
\[
 B(\mathcal T_A)\subseteq\mathcal T_A\cup Z_A
 \subseteq\mathcal T_A\cup\Fix(A^{r_A}).
\]
Moreover, commutation gives
\[
\begin{split}
 A(\Fix(A^{r_A}))&\subseteq\Fix(A^{r_A}),\\
 B(\Fix(A^{r_A}))&\subseteq\Fix(A^{r_A}).
\end{split}
\]
Therefore
\[
 F_A=\mathcal T_A\cup\Fix(A^{r_A})
\]
is finite and forward invariant under both maps.  Construct \(F_B\)
symmetrically, using the eventual \(B\)-cycles met by the forward \(B\)-orbits
of \(\mathcal V(A)\).  We have proved
\begin{equation}\label{term:eq:F-invariance}
 A(F_A)\cup B(F_A)\subseteq F_A,\qquad
 A(F_B)\cup B(F_B)\subseteq F_B.
\end{equation}

\begin{claim}[Turn transport]\label{clm:pl-turn-transport}
\begin{equation}\label{term:eq:turn-transport}
 B(E(A))\subseteq E(A)\cup F_B,\qquad
 A(E(B))\subseteq E(B)\cup F_A.
\end{equation}
Moreover,
\begin{equation}\label{term:eq:R-turn-transport}
 A(E(R))\subseteq E(R)\cup F_A,\qquad
 B(E(R))\subseteq E(R)\cup F_B.
\end{equation}
\end{claim}

\begin{proof}
Suppose \(x\) is an interior turn of \(A\) but not of \(B\).  Flat-freeness
and the absence of a turn make \(B\) locally strictly monotone and
one-to-one at \(x\).  In \(BA\), the inner turn of \(A\) cannot cancel:
both branches approach \(A(x)\) from the same side, on which the germ of
\(B\) has nonzero slope.  Since \(AB=BA\), the map \(AB\) also turns at
\(x\).  Local one-to-one-ness of \(B\) then forces \(B(x)\) to be a turn of
\(A\).  If \(x\) is a turn of both maps, then
\(B(x)\in\mathcal V(B)\subseteq F_B\).  The same inclusion holds when \(x\)
is a domain endpoint, since every domain endpoint lies in \(E(B)\).
This proves the first inclusion in \eqref{term:eq:turn-transport}; the second
is symmetric.  Repeating the same argument for \(AR=RA\) and \(BR=RB\)
gives \eqref{term:eq:R-turn-transport}.
\end{proof}

Set
\begin{equation}\label{term:eq:Markov-K}
 \mathcal K=
 F_A\cup F_B\cup E(A)\cup E(B)\cup E(R)\cup\{p_\theta\}.
\end{equation}
\begin{claim}[Markov-set invariance]\label{clm:pl-Markov-invariance}
\[
 A(\mathcal K)\subseteq\mathcal K,\qquad
 B(\mathcal K)\subseteq\mathcal K.
\]
\end{claim}

\begin{proof}
Equation~\eqref{term:eq:F-invariance} handles \(F_A\) and \(F_B\).  For the
three vertex sets, the required images under \(A\) are
\[
\begin{aligned}
 A(E(A))&=\mathcal V(A)\subseteq\mathcal T_A\subseteq F_A,\\
 A(E(B))&\subseteq E(B)\cup F_A,\\
 A(E(R))&\subseteq E(R)\cup F_A,
\end{aligned}
\]
where the last two lines are
\eqref{term:eq:turn-transport} and
\eqref{term:eq:R-turn-transport}.  Similarly, their images under \(B\) are
\[
\begin{aligned}
 B(E(A))&\subseteq E(A)\cup F_B,\\
 B(E(B))&=\mathcal V(B)\subseteq\mathcal T_B\subseteq F_B,\\
 B(E(R))&\subseteq E(R)\cup F_B.
\end{aligned}
\]
Finally, \(p_\theta\) is fixed by both maps.  Every summand in
\eqref{term:eq:Markov-K} is therefore carried back into \(\mathcal K\).
\end{proof}

Thus \(\mathcal K\) is finite and forward invariant under \(A,B\).  It contains
all turns of both maps.  Moreover, constant nonzero absolute slope excludes
an unrecorded same-direction kink: on two adjacent pieces with the same
monotonicity the signed slopes coincide, and continuity then makes the two
affine formulas identical.  Hence each map is affine and monotone on every
component of $\I\setminus\mathcal K$, and the endpoint images lie in
$\mathcal K$.
Since \(R=AB\), we also have \(R(\mathcal K)\subseteq\mathcal K\).
This is the required common Markov vertex set.
\end{proof}

\subsection{Common substitutions and Cobham periodicity}

\subsubsection{Directed edge substitutions}

Because \(p<c<1\), we have \(p_\theta<1\).  Let \(k_+\) be the least member
of \(\mathcal K\) greater than \(p_\theta\), and put
\[
 e_+=(p_\theta,k_+),\qquad \ell=k_+-p_\theta.
\]
Because \(\mathcal K\) contains every turn of \(A,B,R\), and these maps have
constant absolute slope, their positive selected affine germs extend as
long as the argument remains in \(e_+\).  For \(0\leq t\leq\ell\), choose
\(N\) so large that \(t/\widehat s^{\,N}\) lies in the initial germ.
Every intermediate point
\(p_\theta+t/\widehat s^{\,j}\), \(0\leq j\leq N\), lies in the closure of
\(e_+\); iterating the positive branch of \(R\) in
\eqref{term:eq:Psi} therefore gives
\[
 \Psi(t)=p_\theta+t\qquad(0\leq t\leq\ell).
\]
Thus \(\Psi|_{[0,\ell]}\) parametrizes the directed edge \(e_+\) at unit
speed.  Applying \eqref{term:eq:Psi-equivariance} \(n\) times gives the
corresponding iterate identities.  We now record all the symbolic information
that will be used below in one proposition.

\begin{proposition}[Common directed-edge coding]
\label{term:prop:common-edge-coding}
There are a finite alphabet \(\mathcal A\) of directed
\(\mathcal K\)-edges, a seed \(e_+\in\mathcal A\), two nonerasing
substitutions \(\sigma_A,\sigma_B\) on \(\mathcal A\), and a one-sided word
\(x\in\mathcal A^{\mathbb N}\) with the following properties.
\begin{enumerate}
\item The word \(x\) is the literal unit-speed directed-edge itinerary of
\(\Psi\), beginning with \(e_+\).  Every letter of \(\mathcal A\) occurs
in \(x\).
\item Both substitutions are prolongable and generating on \(e_+\), and
\begin{equation}\label{term:eq:common-fixed-word}
 \sigma_A(x)=x=\sigma_B(x).
\end{equation}
For every \(n\in\mathbb N\), their iterated seed images are the corresponding
prefixes:
\begin{equation}\label{term:eq:substitution-prefixes}
\begin{split}
 \sigma_A^n(e_+)
 &=\text{the itinerary of }\Psi\text{ on }[0,\widehat a^{\,n}\ell],\\
 \sigma_B^n(e_+)
 &=\text{the itinerary of }\Psi\text{ on }[0,\widehat b^{\,n}\ell].
\end{split}
\end{equation}
In particular, the lengths of both prefix families tend to infinity.
\item For either
\[
 (T,s,\sigma_T)=(A,\widehat a,\sigma_A)
 \quad\hbox{or}\quad
 (T,s,\sigma_T)=(B,\widehat b,\sigma_B),
\]
the word \(\sigma_T(e)\), for every \(e\in\mathcal A\), is the complete
literal directed path traced by \(T\) on \(e\), without cancellation.  More
generally, for every \(n\in\mathbb N\), if
\[
 \sigma_T^n(e)=f_1\cdots f_r,\qquad
 d_j=\sum_{i=1}^j |f_i|,\quad d_0=0,
\]
and \(\iota_e:[0,|e|]\to e\) is the unit-speed parametrization in the
displayed orientation, then for every \(1\leq j\leq r\) and
\(d_{j-1}\leq u\leq d_j\),
\begin{equation}\label{term:eq:ordered-edge-realization}
 T^n\!\left(\iota_e(u/s^n)\right)
   =\iota_{f_j}(u-d_{j-1}).
\end{equation}
\item The exact metric laws are
\begin{equation}\label{term:eq:metric-substitution}
 \sum_{f\text{ in }\sigma_A(e)}|f|=\widehat a|e|,
 \qquad
 \sum_{f\text{ in }\sigma_B(e)}|f|=\widehat b|e|.
\end{equation}
\end{enumerate}
\end{proposition}

\begin{proof}
First work on the full alphabet consisting of both orientations of every
\(\mathcal K\)-edge.  For a directed edge \(e\), the map \(A|_e\) is affine
and strict, its endpoints lie in \(\mathcal K\), and its absolute slope is
\(\widehat a\).  Cutting its oriented image at all intervening Markov
vertices gives a nonempty complete literal path; call this word
\(\Sigma_A(e)\).  Define \(\Sigma_B\) in the same way.  Oppositely directed
successive letters, if they arise in a literal description, are retained.

The construction is pointwise, not just endpoint-compatible.  For either
\((T,s,\Sigma_T)=(A,\widehat a,\Sigma_A)\) or
\((T,s,\Sigma_T)=(B,\widehat b,\Sigma_B)\), if
\(\Sigma_T(e)=f_1\cdots f_r\), then the path
\[
 u\longmapsto T\bigl(\iota_e(u/s)\bigr),
 \qquad 0\leq u\leq s|e|,
\]
has speed one, hits the successive Markov vertices at the cumulative times
\(|f_1|,|f_1|+|f_2|,\ldots\), and equals the displayed unit-speed
parametrization of \(f_j\) between the \((j-1)\)-st and \(j\)-th times.
This proves \eqref{term:eq:ordered-edge-realization} for \(n=1\), as well as
\eqref{term:eq:metric-substitution} on the full alphabet.
For \(n=0\), the ordered realization is simply the identity parametrization
of \(e\).

Inductively, replace every letter \(f_j\) in the depth-\(n\) realization by
its one-step literal image.  That block has duration \(s|f_j|\).  Hence its
right endpoint is
\[
 s\bigl(|f_1|+\cdots+|f_j|\bigr)=s d_j,
\]
which is the left endpoint of the following block after the global
rescaling.  Inside the block, writing the local parameter as
\(u=s d_{j-1}+v\), the identity
\[
 T^{n+1}\!\left(\iota_e(u/s^{n+1})\right)
 =T\!\left(T^n\!\left(
   \iota_e((d_{j-1}+v/s)/s^n)\right)\right)
\]
reduces the assertion to the induction hypothesis and the one-step formula.
Thus \eqref{term:eq:ordered-edge-realization} holds for every \(n\).

On the seed edge, the positive selected germs of both \(A\) and \(B\), not
merely the inequalities \(\widehat a,\widehat b>1\), show that
\(\Sigma_A(e_+)\) and \(\Sigma_B(e_+)\) begin with \(e_+\).  Their metric
lengths are respectively \(\widehat a\ell\) and \(\widehat b\ell\), so the
remaining tails are nonempty.  Moreover
\[
 \Psi(\widehat a^{\,n}t)=A^n(\Psi(t)),\qquad
 \Psi(\widehat b^{\,n}t)=B^n(\Psi(t))
 \quad(0\leq t\leq\ell).
\]
The ordered realization therefore identifies the two iterated seed words
with the prefixes in \eqref{term:eq:substitution-prefixes}.

We justify explicitly that the ray has a literal \(\mathcal K\)-edge
itinerary.  On a compact parameter interval, choose one sufficiently large
\(N\) in \eqref{term:eq:Psi}, so that \(\Psi\) there is the affine initial
segment followed by \(R^N\).  A breakpoint of this composite occurs when
some intermediate iterate meets \(E(R)\).  Since
\(E(R)\subseteq\mathcal K\) and \(R(\mathcal K)\subseteq\mathcal K\), its
final \(\Psi\)-value belongs to \(\mathcal K\).  Refine also at the finitely
many times at which an affine piece crosses a vertex of \(\mathcal K\).
Apart from a possibly truncated final piece at the chosen compact endpoint,
the resulting consecutive pieces are precisely complete directed
\(\mathcal K\)-edges.  These subdivisions are compatible as the compact
parameter interval grows.  The finite alphabet has a positive minimum edge
length, so completed edges cannot accumulate in bounded parameter time;
hence they form a well-defined one-sided word \(x\), beginning with \(e_+\).

The substitution blocks align pointwise.  If a ray letter \(e\) occupies
the parameter interval \([r,r+|e|]\), its complete \(A\)-image occupies
\([\widehat a r,\widehat a(r+|e|)]\) and parametrizes
\(u\mapsto A(\Psi(u/\widehat a))\).  Adjacent blocks meet at the image under
\(A\) of their common endpoint, so their literal concatenation has no gap; a
possible immediate reversal is retained rather than cancelled.  By
\(\Psi(\widehat a t)=A(\Psi(t))\), the
concatenated path is exactly \(\Psi\).  Literal-itinerary uniqueness gives
\(\Sigma_A(x)=x\), and the same argument gives \(\Sigma_B(x)=x\).

Let \(\mathcal A\) be the set of directed letters occurring in \(x\).  If
\(e\) occurs in \(x\), every letter of \(\Sigma_A(e)\) occurs in
\(\Sigma_A(x)=x\), and similarly for \(B\).  Thus \(\mathcal A\) is
invariant under both substitutions, and restriction gives
\(\sigma_A,\sigma_B\) while preserving the pointwise and metric identities.
The seed belongs to \(\mathcal A\); the positive-germ observation makes both
restrictions prolongable on it.  Finally the metric lengths of the seed
prefixes are \(\widehat a^{\,n}\ell\) and
\(\widehat b^{\,n}\ell\).  Since the finite alphabet has a positive minimum
and finite maximum edge length, their ordinary word lengths tend to
infinity.  Every letter occurring in \(x\) consequently occurs in some
iterated seed prefix, so \(e_+\) is generating for both substitutions.
This proves all assertions.
\end{proof}

\subsubsection{Perron growth}

Use the standard incidence convention
$(M_A)_{e,f}=|\sigma_A(f)|_e$, and similarly for $B$: columns are source
letters and rows count output letters.  Assign to each directed geometric
edge its positive Euclidean length, and let $\boldsymbol\ell$ be the
resulting strictly positive column vector.  Since the image path of an edge
has Euclidean length equal to the absolute slope times the length of that
edge,
\begin{equation}\label{term:eq:length-eigenvectors}
 M_A^{\mathsf T}\boldsymbol\ell=\widehat a\boldsymbol\ell,\qquad
 M_B^{\mathsf T}\boldsymbol\ell=\widehat b\boldsymbol\ell.
\end{equation}

Recall that a \emph{Perron number} is a real algebraic integer
\(\lambda>1\) all of whose algebraic conjugates other than \(\lambda\) have
modulus strictly less than \(\lambda\).

\begin{lemma}[A weak Perron eigenvalue has a Perron power]
\label{term:lem:perron-power}
Let \(M\) be a finite nonnegative integral matrix.  If
\[
 Mv=\lambda v
\]
for a strictly positive vector \(v\) and a real number \(\lambda>1\), then
\(\rho(M)=\lambda\), and there is a positive integer \(u\) such that
\(\lambda^u\) is a Perron number.  Moreover \(M^u\) has spectral radius
\(\lambda^u\).
\end{lemma}

\begin{proof}
The Collatz--Wielandt inequalities applied to the strictly positive vector
\(v\) give \(\rho(M)=\lambda\).  Since \(M\) is integral, \(\lambda\) is an
algebraic integer; its minimal polynomial divides the characteristic
polynomial of \(M\), so every algebraic conjugate of \(\lambda\) is an
eigenvalue of \(M\) and has modulus at most \(\lambda\).

Put \(M\) in Frobenius block form.  If an eigenvalue \(z\) of \(M\) has
\(|z|=\lambda\), then it belongs to an irreducible diagonal block whose
spectral radius is \(\lambda\).  The periodic form of the Perron--Frobenius
theorem gives \(z=\lambda\zeta\), where \(\zeta\) is a root of unity.  Choose
\(u\) divisible by the orders of all the finitely many such \(\zeta\).

Let \(w\) be a conjugate of \(\lambda^u\).  An embedding of
\(\mathbb Q(\lambda^u)\) producing \(w\) extends to an embedding of
\(\mathbb Q(\lambda)\), so \(w=z^u\) for a conjugate \(z\) of \(\lambda\).
If \(|z|<\lambda\), then \(|w|<\lambda^u\).  If \(|z|=\lambda\), the
choice of \(u\) gives \(w=z^u=\lambda^u\).  Thus every distinct conjugate of
\(\lambda^u\) has modulus strictly less than \(\lambda^u\), which is exactly
the Perron property.  Finally
\(\rho(M^u)=\rho(M)^u=\lambda^u\).
\end{proof}

Apply Lemma~\ref{term:lem:perron-power} to
$M_A^{\mathsf T},M_B^{\mathsf T}$.  Transposition does not change the
spectrum.  There are positive integers $u,v$ such that
$\widehat a^{\,u},\widehat b^{\,v}$ are Perron numbers and are the spectral
radii of $M_A^u,M_B^v$, respectively.  The incidence matrix of a substitution
power is the corresponding matrix power.
Moreover
\[
 \sigma_A^u(x)=x=\sigma_B^v(x),
\]
the occurring alphabet is unchanged, and the substitutions remain
nonerasing and prolongable on \(e_+\).  Their iterated prefixes form
subsequences of the prefixes in \eqref{term:eq:substitution-prefixes}.  The
positive minimum and maximum edge lengths show that
\(\lvert\sigma_A^n(e_+)\rvert\) is comparable to \(\widehat a^{\,n}\), and
similarly for \(B\), so both prefix lengths tend to infinity.  The same is
therefore true along the subsequences of exponents divisible by \(u\) and
\(v\).  Since those subsequence images are prefixes of \(x\) of unbounded
length, every retained letter occurs in an iterated seed image of each
powered substitution.  Thus the powered presentations are generating at
the actual head seed.  Finally \(\widehat a^{\,u}\) and
\(\widehat b^{\,v}\) remain multiplicatively independent, since a relation
between positive powers of them would be a relation between positive powers
of \(\widehat a,\widehat b\).

\subsubsection{Durand's theorem and periodicity}

We use Durand's substitution form of Cobham's theorem
\cite[Definition 2 and Theorem 1]{Durand}.

\medskip
\noindent\textbf{Durand--Cobham theorem.}
\emph{Let \(\alpha,\beta>1\) be multiplicatively independent Perron
numbers.  Suppose a one-sided sequence over a finite alphabet is both
\(\alpha\)-substitutive and \(\beta\)-substitutive: more precisely, it is a
letter-to-letter coding of a fixed point of a substitution prolongable on a
generating letter whose occurring-letter incidence matrix has dominant
eigenvalue \(\alpha\), and it has an analogous presentation with dominant
eigenvalue \(\beta\).  Then the sequence is ultimately periodic.}
\medskip

Set \(\tau_A=\sigma_A^u\) and \(\tau_B=\sigma_B^v\).  We verify the
Definition~2 boundary explicitly.  Both substitutions are nonerasing and
prolongable on the actual first letter \(e_+\).  Their canonical generated
fixed points are \(x\), because \(\tau_A^n(e_+)=\sigma_A^{un}(e_+)\) and
\(\tau_B^n(e_+)=\sigma_B^{vn}(e_+)\) are cofinal families of prefixes of
\(x\), and their lengths tend to infinity.  Every letter in the retained
alphabet occurs in \(x\), and the preceding paragraph shows that \(e_+\) is
generating for both powered substitutions.  The incidence matrices are
\(M_A^u,M_B^v\); their dominant eigenvalues are respectively
\(\widehat a^u,\widehat b^v\), which are Perron numbers by
Lemma~\ref{term:lem:perron-power}.  Finally the output coding in both
presentations is the identity letter coding.  These are exactly the two
substitutive presentations required in Durand's Definition~2.  Their
dominant eigenvalues are multiplicatively independent, so Durand's
Theorem~1 makes the directed-edge word \(x\) ultimately periodic.

\begin{proposition}[Least-period periodic return path]
\label{term:prop:periodic-return-path}
There are a continuous surjection \(\psi:\mathbb R\to\I\) and a number
\(P_0>0\) such that:
\begin{enumerate}
\item \(\psi\) is finite-piece affine with slopes \(+1\) or \(-1\) on every
compact parameter interval;
\item its full period group is \(P_0\mathbb Z\);
\item \(\widehat a,\widehat b\) are natural numbers greater than one; and
\item for every \(t\in\mathbb R\),
\begin{equation}\label{term:eq:psi-returns}
 \psi(\widehat a t)=A(\psi(t)),\qquad
 \psi(\widehat b t)=B(\psi(t)).
\end{equation}
\end{enumerate}
Thus \(\psi\) descends to an onto unit-speed polygonal circle path
\begin{equation}\label{term:eq:psi-circle}
 \psi\colon\mathbb R/P_0\mathbb Z\longrightarrow\I.
\end{equation}
\end{proposition}

\begin{proof}
Choose integers \(N\geq0\) and \(q\geq1\) such that
\(x_{j+q}=x_j\) for every \(j\geq N\).  Put
\[
 \tau_n=\sum_{0\leq j<n}|x_j|,
 \qquad
 P=\tau_{N+q}-\tau_N
   =\sum_{N\leq j<N+q}|x_j|>0.
\]
Periodicity of the letters gives
\(\tau_{n+q}=\tau_n+P\) for every \(n\geq N\).  The positive minimum length
of a directed \(\mathcal K\)-edge gives \(\tau_n\to\infty\).  Hence every
\(t\geq\tau_N\) lies in some \([\tau_n,\tau_{n+1}]\) with \(n\geq N\).
On such an interval, the literal itinerary formula says that
\(\Psi(t)\) is the point at offset \(t-\tau_n\) on the directed edge
\(x_n\).  The same description at \(n+q\) therefore gives
\begin{equation}\label{term:eq:eventual-metric-period}
 \Psi(t+P)=\Psi(t)\qquad(t\geq\tau_N).
\end{equation}
This includes vertex times, because each directed-edge parametrization
includes both endpoints.

\smallskip
\noindent\emph{Periodic extension and local unit speed.}
For \(t\in\mathbb R\), choose any integer \(m\geq0\) large enough that
\(t+mP\geq\tau_N\), and define
\begin{equation}\label{term:eq:periodic-tail-extension}
 \psi(t)=\Psi(t+mP).
\end{equation}
If \(m'\geq m\) is another admissible choice, repeated use of
\eqref{term:eq:eventual-metric-period}, starting at the already late
parameter \(t+mP\), gives the same value.  Thus the definition is
independent of \(m\).  Choosing one \(m\) which works on a neighborhood of a
given \(t\) identifies \(\psi\) there with a translate of \(\Psi\); hence
\(\psi\) is continuous.  The same argument with one \(m\) working on a
given compact interval transfers a finite unit-speed subdivision from
\(\Psi\) to \(\psi\).  Taking the same sufficiently large shift in
\eqref{term:eq:periodic-tail-extension} at \(t\) and \(t+P\) proves
\(\psi(t+P)=\psi(t)\) for every real \(t\).  Finally \(\psi=\Psi\) on the
late tail.  A unit-slope affine subinterval of that tail shows that
\(\psi\) is nonconstant.

\smallskip
\noindent\emph{The periodic tail is onto.}
Let
\[
 J=\Psi([\tau_N,\infty))=\psi(\mathbb R).
\]
The period just proved shows that \(J=\Psi([\tau_N,\tau_N+P])\), so \(J\)
is a compact interval.  It has nonempty relative interior because the late
ray has a nondegenerate unit-slope affine piece.  Since
\(\widehat s=\widehat a\widehat b>1\), for \(t\geq\tau_N\) we have
\(\widehat s t\geq\tau_N\), and the product return gives
\[
 R(\Psi(t))=\Psi(\widehat s t)\in J.
\]
Thus \(R(J)\subseteq J\).  If \(J\ne\I\), choose a nonempty relatively open
interval \(U\) with \(U\subseteq\operatorname{int}_{\I}J\).  Exactness of
\(R\) gives an \(n\) for which \(R^n(U)=\I\), while forward invariance gives
\(R^n(U)\subseteq J\), a contradiction.  Hence \(J=\I\), proving that
\(\psi\) is onto.

\smallskip
\noindent\emph{Least period and full period group.}
Set
\[
 \operatorname{Per}(\psi)
 =\{r\in\mathbb R:\psi(t+r)=\psi(t)\text{ for every }t\in\mathbb R\}.
\]
This is an additive subgroup containing \(P\).  It is closed: if periods
\(r_i\) converge to \(r\), continuity gives
\(\psi(t+r)=\lim_i\psi(t+r_i)=\psi(t)\) for every \(t\).
It is not dense, for otherwise continuity would make every real number a
period and \(\psi\) would be constant.  We spell out the elementary subgroup
classification used here.  If every neighborhood of zero contained a
nonzero period, we could choose periods \(h_n\) with
\(0<|h_n|<1/n\).  For any \(x\in\mathbb R\), choose an integer \(k_n\)
nearest to \(x/h_n\).  Then
\[
 |x-k_nh_n|\leq |h_n|/2\longrightarrow0,
\]
so the period group would be dense, a contradiction.  Hence there is an
\(\varepsilon>0\) such that zero is its only member in
\((-\varepsilon,\varepsilon)\).

The set of positive periods is nonempty because it contains \(P\).  Let
\[
 P_0=\inf\bigl(\operatorname{Per}(\psi)\cap(0,\infty)\bigr).
\]
The preceding separation gives \(P_0\geq\varepsilon>0\).  By the defining
property of the infimum, there are positive periods \(r_n\) with
\(P_0\leq r_n<P_0+1/n\).  Thus \(r_n\to P_0\), and closedness makes \(P_0\)
a period.  It is consequently the least positive period.  For any period
\(r\), put \(k=\lfloor r/P_0\rfloor\).  The remainder
\(r-kP_0\) is again a period and lies in \([0,P_0)\), so leastness forces it
to be zero.  Conversely every integral multiple of \(P_0\) is a period.
Therefore
\[
 \operatorname{Per}(\psi)=P_0\mathbb Z.
\]

\smallskip
\noindent\emph{Integral scales.}
Because \(\psi\) agrees with \(\Psi\) sufficiently far out, the ray returns
give cutoffs \(T_A,T_B\) such that
\[
 \psi(\widehat a t)=A(\psi(t))\quad(t\geq T_A),
 \qquad
 \psi(\widehat b t)=B(\psi(t))\quad(t\geq T_B).
\]
For \(t\geq T_A\), periodicity gives
\[
\begin{split}
 \psi\bigl(\widehat a t+\widehat aP_0\bigr)
 &=\psi\bigl(\widehat a(t+P_0)\bigr)\\
 &=A(\psi(t+P_0))
 =A(\psi(t))
 =\psi(\widehat a t).
\end{split}
\]
Thus translation by \(\widehat aP_0\) fixes \(\psi\) at every sufficiently
large argument.  Given arbitrary \(y\), choose an integer \(m\geq0\) with
\(y+mP_0\) in that late range.  Translating both arguments by \(mP_0\)
shows \(\psi(y+\widehat aP_0)=\psi(y)\).  Hence
\(\widehat aP_0\in\operatorname{Per}(\psi)=P_0\mathbb Z\).  Since \(P_0>0\),
there is an integer \(k_A>1\) with \(\widehat a=k_A\).  The identical
argument gives an integer \(k_B>1\) with \(\widehat b=k_B\).  Therefore
\begin{equation}\label{term:eq:integer-slopes}
 \widehat a,\widehat b\in\mathbb N.
\end{equation}

\smallskip
\noindent\emph{Global return identities.}
Let \(t\in\mathbb R\), and choose \(m\geq0\) with \(t+mP_0\geq T_A\).
Since \(\widehat a\) is an integer,
\[
\begin{split}
 \psi(\widehat a t)
 &=\psi\bigl(\widehat a t+\widehat a mP_0\bigr)\\
 &=\psi\bigl(\widehat a(t+mP_0)\bigr)
 =A(\psi(t+mP_0))
 =A(\psi(t)).
\end{split}
\]
The same calculation for \(B\) proves \eqref{term:eq:psi-returns} on all of
\(\mathbb R\), completing the proof.
\end{proof}

\subsection{Original-phase affine lifts and the common fixed point}

\subsubsection{Cross-period globalization}

\begin{lemma}[Cross-period globalization]\label{term:lem:cross-period}
Let \(u,v\colon\mathbb R\to\I\) be continuous periodic surjections.  Suppose
that on every compact parameter interval each path has a finite affine
subdivision and every affine slope is \(+1\) or \(-1\).  If, for some
\(r>1\) and some interval map \(T\),
\[
 u(rt)=T(u(t)),\qquad v(rt)=T(v(t)).
\]
Then \(u\) and \(v\), up to a translation and possibly a reflection of the
parameter, are the same path.  In particular, they have the same full period
group and the same least positive period.
\end{lemma}

\begin{proof}
Choose positive periods \(P_u,P_v\).  Each path has only finitely many
vertices in a fundamental interval, and periodicity therefore makes the set
of all their vertex images finite.  Choose \(y\in(0,1)\) outside the vertex
images of both paths.
Surjectivity supplies preimages \(t_u,t_v\).  They lie in interiors of
unit-slope affine cells, so, after shrinking to a common nonempty interval,
there are
$\varepsilon\in\{1,-1\}$ and $\delta\in\mathbb R$ such that
\begin{equation}\label{term:eq:local-path-relation}
 v(t)=u(\varepsilon t+\delta)
\end{equation}
on a nonempty open interval $J$.

Iterating the two dilation identities shows that, on $J_n=r^nJ$,
\begin{equation}\label{term:eq:expanded-path-relation}
 v(t)=u(\varepsilon t+r^n\delta).
\end{equation}
Indeed, write $t=r^ns$ with $s\in J$ and apply $T^n$ to
\eqref{term:eq:local-path-relation}.  The lengths of the intervals $J_n$
tend to infinity.

Let $m_n$ be the midpoint of $J_n$.  Reduce
\[
 \theta_n=m_n\pmod {P_v},\qquad
 \eta_n=\varepsilon m_n+r^n\delta\pmod {P_u}
\]
to fixed compact fundamental intervals.  Pass to a subsequence on which
$\theta_n\to\theta$ and $\eta_n\to\eta$.  For any fixed
$h\in\mathbb R$, the point $m_n+h$ belongs to $J_n$ for all sufficiently
large $n$.  Periodicity and \eqref{term:eq:expanded-path-relation} give
\[
 v(\theta_n+h)=u(\eta_n+\varepsilon h).
\]
Taking the limit yields
\[
 v(\theta+h)=u(\eta+\varepsilon h)
 \qquad(h\in\mathbb R).
\]
Thus, with $\delta_0=\eta-\varepsilon\theta$,
\[
 v(t)=u(\varepsilon t+\delta_0)\qquad(t\in\mathbb R).
\]
A number is a period of one side exactly when its signed counterpart is a
period of the other.  Period groups are closed under negation, so the two
full period groups coincide.  In particular their least positive periods
coincide.
\end{proof}

\subsubsection{Affine circle lifts of the original maps}

\begin{lemma}[Rescaling an original map]\label{term:lem:rescaled-original-path}
Let \(\varphi:\mathbb R\to\I\) be an onto periodic path which is
finite-piece affine with slopes \(+1\) or \(-1\) on every compact interval.
Let \(L\) be a positive integer, let \(P\) be a finite-piece PL map of
constant absolute slope \(C>1\), and suppose
\[
 A=P^L,\qquad a=C^L,\qquad
 \varphi(at)=A(\varphi(t))\quad(t\in\mathbb R).
\]
Then
\[
 v(t)=P(\varphi(t/C))
\]
is a continuous onto periodic path, finite-piece affine with slopes
\(+1\) or \(-1\) on every compact interval.  For every
\(P_0\in\mathbb R\) that is a period of \(\varphi\), the number \(CP_0\) is
a period of \(v\), and
\[
 v(at)=A(v(t))\qquad(t\in\mathbb R).
\]
\end{lemma}

\begin{proof}
The return identity and onto-ness of \(\varphi\) first imply that \(A\) is
onto.  Indeed, if \(\varphi(t)=y\), then
\[
 A\bigl(\varphi(t/a)\bigr)=\varphi(t)=y.
\]
Since \(L\geq1\) and \(A=P^L=P\circ P^{L-1}\), onto-ness of \(A\) implies
onto-ness of \(P\).  Given \(y\), choose \(z\) with \(P(z)=y\), then choose
\(s\) with \(\varphi(s)=z\); now \(v(Cs)=y\).  Thus \(v\) is onto.
Continuity is immediate.

Fix a compact parameter interval \(K\).  Subdivide it so that
\(t\mapsto\varphi(t/C)\) is affine with slope \(+1/C\) or \(-1/C\) on each
cell.  Refine once more at the preimages of the finitely many vertices of
\(P\).  On an affine cell the first map is strict, so each vertex of \(P\)
has at most one preimage; the refinement is still finite.  On every refined
cell the signed slope of \(P\) is \(+C\) or \(-C\), and consequently the
signed slope of the composite \(v\) is \(+1\) or \(-1\).

Let \(P_0\) be any period of \(\varphi\).  Then
\[
 v(t+CP_0)
 =P\bigl(\varphi(t/C+P_0)\bigr)
 =P\bigl(\varphi(t/C)\bigr)
 =v(t).
\]
Finally, \(P\) commutes with its iterate \(A=P^L\), and therefore
\[
\begin{split}
 v(at)
 &=P\bigl(\varphi(at/C)\bigr)
  =P\bigl(\varphi(a(t/C))\bigr)\\
 &=P\bigl(A(\varphi(t/C))\bigr)
  =A\bigl(P(\varphi(t/C))\bigr)
  =A(v(t)).
\end{split}
\]
\end{proof}

Recall that \(P,Q\) have constant absolute slopes \(C,D>1\), with
\begin{equation}\label{term:eq:root-slopes}
 C^L=\widehat a,\qquad D^L=\widehat b.
\end{equation}
Define
\[
 \psi_P(t)=P(\psi(t/C)),\qquad
 \psi_Q(t)=Q(\psi(t/D)).
\]
Lemma~\ref{term:lem:rescaled-original-path}, applied once to \(P,A,C\) and
once to \(Q,B,D\), shows that these are onto periodic locally unit-speed
polygonal paths, that \(CP_0,DP_0\) are respective periods, and that they
have the required pairwise self-similarities: \(\psi_P\) has the same
\(A\)-self-similarity as \(\psi\), and \(\psi_Q\) has the same
\(B\)-self-similarity as \(\psi\).

Apply Lemma~\ref{term:lem:cross-period} to \(\psi_P,\psi\), with
\(r=\widehat a\) and \(T=A\).  There are
\(\varepsilon_P\in\{1,-1\}\) and \(\delta_P\in\mathbb R\) such that
\[
 \psi_P(t)=\psi(\varepsilon_Pt+\delta_P)\qquad(t\in\mathbb R),
\]
and the two paths have the same full period group \(P_0\mathbb Z\).
Since \(CP_0\) is a period of \(\psi_P\), it is therefore an integral
multiple of \(P_0\).  Cancellation of \(P_0>0\) shows that \(C\) is a
positive integer; since \(C>1\), it is at least two.  Substituting \(t=Cx\)
in the global path relation gives an affine circle lift
\begin{equation}\label{term:eq:P-lift}
 P(\psi(x))=\psi(\mathcal P(x)),\qquad
 \mathcal P(x)=kx+\alpha\pmod{P_0},
\end{equation}
where \(k=\varepsilon_PC\in\mathbb Z\), \(|k|\geq2\), and we choose the
real representative \(\alpha=\delta_P\) of the translation class.  The same
argument with \(\psi_Q,\psi\) gives
\begin{equation}\label{term:eq:Q-lift}
 Q(\psi(x))=\psi(\mathcal Q(x)),\qquad
 \mathcal Q(x)=\ell x+\beta\pmod{P_0},
\end{equation}
where \(\ell\in\mathbb Z\), \(|\ell|\geq2\), and
\(\beta\in\mathbb R\) is a fixed real representative of the corresponding
translation class.

Physical commutation makes these affine lifts commute exactly.  With the
displayed real representatives fixed, the real affine representatives of
\(\mathcal P\mathcal Q\) and \(\mathcal Q\mathcal P\) have common linear
coefficient \(k\ell\), and differ by the constant
\[
 \Delta=(k\beta+\alpha)-(\ell\alpha+\beta).
\]
Equations \eqref{term:eq:P-lift}--\eqref{term:eq:Q-lift} and \(PQ=QP\)
show that their images under \(\psi\) agree.  Since \(k\ell\ne0\), either
affine composition is onto on the parameter circle.  Given a real \(t\),
choose a parameter whose image under the first composition is congruent to
\(t\) modulo \(P_0\).  Equality of the two physical images and
\(P_0\)-periodicity give \(\psi(t)=\psi(t-\Delta)\).  Replacing \(t\) by
\(t+\Delta\) yields
\[
 \psi(t+\Delta)=\psi(t)\qquad\text{for every }t.
\]
The full period group gives \(\Delta\in P_0\mathbb Z\), which is precisely
\(\mathcal P\mathcal Q=\mathcal Q\mathcal P\) on the parameter circle.

\subsubsection{The unique reflection}

\begin{lemma}[Globalization and uniqueness of the reflection]
\label{term:lem:global-reflection}
Let \(P_0>0\), and let \(\psi:\mathbb R\to\I\) be a continuous surjection
which is finite-piece affine with slopes \(+1\) or \(-1\) on every compact
parameter interval and whose full period group is \(P_0\mathbb Z\).  Suppose
that, for an interval map \(T\) and a number \(q>1\),
\[
 \psi(qt)=T(\psi(t))\qquad(t\in\mathbb R).
\]
Then there is a unique reflection of the parameter circle
\[
 \rho(t)=-t+\gamma\pmod{P_0}
\]
such that \(\psi\circ\rho=\psi\).
\end{lemma}

\begin{proof}
Choose \(t_0\) with \(\psi(t_0)=1\), using onto-ness.  Take a finite affine
subdivision of \([t_0-1,t_0+1]\) and refine it to contain \(t_0\).  On the
cell immediately to the left of \(t_0\), the slope must be
\(+1\): slope \(-1\) would force values greater than \(1\).  On a
cell immediately to the right, the slope must
be \(-1\) for the same reason.  Hence there is \(d>0\) such that
\[
 \psi(t_0-r)=1-r=\psi(t_0+r)
 \qquad(0\leq r\leq d).
\]
With
\[
 J_0=[t_0,t_0+d],\qquad \gamma_0=2t_0,
\]
we therefore have
\begin{equation}\label{term:eq:local-reflection}
 \psi(t)=\psi(-t+\gamma_0)\qquad(t\in J_0).
\end{equation}

Iterating the displayed self-similarity gives, for every
\(n\geq0\),
\[
 \psi(q^nt)=T^n(\psi(t)).
\]
Apply this to both sides of \eqref{term:eq:local-reflection}.  On
\[
 J_n=q^nJ_0,\qquad
 \gamma_n=q^n\gamma_0,
\]
it gives
\begin{equation}\label{term:eq:expanded-reflection}
 \psi(u)=\psi(-u+\gamma_n)\qquad(u\in J_n).
\end{equation}
Choose \(n\) so large that the length of \(J_n\) exceeds \(P_0\).  Given any
\(t\in\mathbb R\), choose \(m\in\mathbb Z\) with \(t+mP_0\in J_n\).
Periodicity and \eqref{term:eq:expanded-reflection} yield
\[
\begin{split}
 \psi(t)
 &=\psi(t+mP_0)
 =\psi(-(t+mP_0)+\gamma_n)\\
 &=\psi(-t+\gamma_n).
\end{split}
\]
Thus \(\rho(t)=-t+\gamma_n\) is a global reflection on
\(\mathbb R/P_0\mathbb Z\).

For uniqueness, suppose \(\rho_\gamma,\rho_\delta\) are two global reflection
symmetries.  Their composition is translation by \(\delta-\gamma\), so
\(\delta-\gamma\) is a period of \(\psi\).  The assumed full period group of
this path is \(P_0\mathbb Z\), so \(\delta-\gamma\in P_0\mathbb Z\), and the two
reflections are identical on the parameter circle.
\end{proof}

Choose a real representative \(\gamma\) of the translation class of the
unique reflection, and let
\begin{equation}\label{term:eq:reflection}
 \rho(t)=-t+\gamma\pmod{P_0},\qquad \psi\circ\rho=\psi
\end{equation}
be the unique reflection supplied by
Lemma~\ref{term:lem:global-reflection}, applied with
\(q=\widehat a\) and \(T=A\).

We now prove that the lift \(\mathcal P\) preserves the \(\rho\)-orbits.
First,
\[
\begin{split}
 \psi(\mathcal P\rho(t))
 &=P(\psi(\rho(t)))\\
 &=P(\psi(t))
 =\psi(\mathcal P(t)).
\end{split}
\]
Working on the parameter circle, write
\(\mathcal P(t)=kt+\alpha\) and \(\rho(t)=-t+\gamma\) modulo \(P_0\).
If \(y=\mathcal P(t)\) as a circle class, then, without choosing an inverse
branch,
\[
 \mathcal P\rho(t)=-y+k\gamma+2\alpha\pmod{P_0}.
\]
Thus
\[
 \rho_P(y)=-y+k\gamma+2\alpha\pmod{P_0}
\]
is another global reflection symmetry of \(\psi\): explicitly,
\(\mathcal P\rho(t)=\rho_P(\mathcal P(t))\), and every \(y\) equals
\(\mathcal P(t)\) for some \(t\), so \(\psi(\rho_P(y))=\psi(y)\).
Uniqueness gives \(\rho_P=\rho\), and therefore
\[
 \mathcal P\rho=\rho\mathcal P.
\]
The same argument gives \(\mathcal Q\rho=\rho\mathcal Q\).

\subsubsection{A common fixed reflection orbit}

Normalize the parameter circle in \eqref{term:eq:psi-circle} by the affine coordinate
\(t\mapsto P_0t+\gamma/2\) from \(\mathbb R/\mathbb Z\) to
\(\mathbb R/P_0\mathbb Z\).  In this coordinate \(\rho(t)=-t\).  Conjugate
\(\mathcal P,\mathcal Q\) by the same coordinate change and continue to use
the symbols \(\psi,\mathcal P,\mathcal Q\) for the normalized objects.  The
physical lift identities \eqref{term:eq:P-lift}--\eqref{term:eq:Q-lift}, the
relations \(\mathcal P\rho=\rho\mathcal P\) and
\(\mathcal Q\rho=\rho\mathcal Q\), and exact commutation of the two lifts are
all preserved by conjugacy; their degrees remain \(k,\ell\).  Thus
\(\psi(-t)=\psi(t)\).  Choose representatives
\(\alpha,\beta\in[0,1)\) for the new translation classes and write
\[
 \mathcal P(t)=kt+\alpha,\qquad
 \mathcal Q(t)=\ell t+\beta\pmod1.
\]
The identity \(\mathcal P\rho=\rho\mathcal P\), evaluated at any \(t\),
gives
\[
 -kt+\alpha=-kt-\alpha\pmod1,
\]
so \(2\alpha=0\pmod1\).  The corresponding calculation for \(\mathcal Q\)
gives \(2\beta=0\pmod1\).  Hence
\begin{equation}\label{term:eq:half-translations}
 \alpha,\beta\in\{0,\tfrac12\},
\end{equation}
while exact commutation of the two affine lifts gives
\begin{equation}\label{term:eq:affine-commutation}
 (k-1)\beta=(\ell-1)\alpha\pmod1.
\end{equation}
The following exhaustive parity check produces an orbit
\(\{t,-t\}\) preserved by both affine maps.
\begin{enumerate}
 \item[(i)] If \((\alpha,\beta)=(0,0)\), take \(t=0\).
 \item[(ii)] If \((\alpha,\beta)=(\tfrac12,\tfrac12)\), then
 \eqref{term:eq:affine-commutation} says that \(k,\ell\) have the same
 parity.  Take \(t=\tfrac12\) if both are even and \(t=\tfrac14\) if both
 are odd.
 \item[(iii)] If \((\alpha,\beta)=(0,\tfrac12)\), then
 \eqref{term:eq:affine-commutation} says that \(k\) is odd.  Take
 \(t=\tfrac12\) if \(\ell\) is even and \(t=\tfrac14\) if \(\ell\) is odd.
 \item[(iv)] If \((\alpha,\beta)=(\tfrac12,0)\), then
 \eqref{term:eq:affine-commutation} says that \(\ell\) is odd.  Take
 \(t=\tfrac12\) if \(k\) is even and \(t=\tfrac14\) if \(k\) is odd.
\end{enumerate}
This calculation uses only parity and therefore includes negative degrees.
For the resulting \(t\), equations \eqref{term:eq:P-lift} and
\eqref{term:eq:Q-lift}, together with \(\psi(-t)=\psi(t)\), give
\[
 P(\psi(t))=\psi(t)=Q(\psi(t)).
\]
This contradicts the standing assumption that \(P,Q\) have no common fixed
point and completes the proof of Proposition~\ref{prop:PL-core}.
\hfill\(\square\)\par


\clearpage
\section{Reduction to the piecewise-linear core}

The following two published results supply the entropy and constant-slope
inputs for the reduction.  All arguments needed to make the semiconjugacy
simultaneous for $P$ and $Q$ are included below.

The entropy input is Theorem~1 of Bogatyi--Frolkina
\cite[Theorem~1]{BogatyiFrolkina}.
\begin{theorem}[Bogatyi--Frolkina]
\label{pm:thm:BF}
If two commuting continuous selfmaps $F,G$ of a compact interval have no
common fixed point, then
\[
 h_2(F\circ G)=h_2(G\circ F)\geq2,
\]
where Bogatyi and Frolkina use logarithms to base two in the definition of
entropy.  In the natural-logarithm convention this says
\(h_{\rm top}(F\circ G)\geq2\log2>0\).
\end{theorem}

The constant-slope input combines Theorem~A and Lemmas~1.1, 3.1, and 4.2
of Alsed\`a--Misiurewicz
\cite[Theorem~A and Lemmas~1.1, 3.1, 4.2]{AlsedaMisiurewicz}.
\begin{theorem}[Alsed\`a--Misiurewicz]
\label{pm:thm:constant-slope}
Let $H$ be a continuous finite-piece piecewise-monotone interval map with
$h_{\rm top}(H)>0$, and put
\[
 \lambda=\exp(h_{\rm top}(H))>1.
\]
Then there are a continuous nondecreasing surjection
$\psi:\I\to\I$ and a continuous finite-piece piecewise-linear map
$h:\I\to\I$, whose nonzero slopes have absolute value $\lambda$, such that
\[
 \psi H=h\psi.
\]
For all $x<y$ contained in one monotonicity piece of $H$, the
constant-slope equation gives
\begin{equation}\label{pm:eq:semiconjugacy-scaling}
 |\psi(H(y))-\psi(H(x))|
 =\lambda\,|\psi(y)-\psi(x)|.
\end{equation}
In addition, $\psi$ may be chosen $H$-compatible in the sense of
Alsed\`a--Misiurewicz.
\end{theorem}

For clarity about the cited input, Lemma~3.1 of the cited paper first
semiconjugates a positive-entropy map with possibly constant laps to a
piecewise strictly monotone map of the same entropy, through an
$H$-compatible increasing map.  Theorem~A supplies the compatible
constant-slope semiconjugacy with slope
\(\exp(h_{\rm top}(H))\).  Lemma~1.1 preserves the orientation of each
noncollapsed lap.  The paper's equation~(2), used in Lemma~4.2, supplies
\eqref{pm:eq:semiconjugacy-scaling}; this scaling identity comes from the
constant-slope factor.  Separately, compatibility ensures that a lap
collapsed by the first semiconjugacy has collapsed image as well.  Thus
constant monotonicity pieces are included without attributing branch
scaling to compatibility.

\subsection{Minimal invariant core}

Suppose, toward a contradiction, that $P,Q$ commute and have no common fixed
point.  Let $\mathscr J$ be the family of nonempty compact intervals $J$
satisfying
\[
 P(J)\subseteq J,\qquad Q(J)\subseteq J,
\]
ordered by inclusion.  Every chain has a nonempty compact intersection, and
that intersection again belongs to $\mathscr J$: for a chain, every finite
intersection is its least member, so compactness gives a nonempty total
intersection; and if $x$ lies in every member, then so do $P(x)$ and $Q(x)$.
Thus every chain has a lower bound.  By Zorn's lemma there is an
inclusion-minimal member.  It is nondegenerate, because an invariant
singleton is a common fixed point.  Restrict to this interval and rescale it
affinely to $\I$.  We continue to write $P,Q$ for the restricted maps.

Minimality implies
\begin{equation}\label{pm:eq:onto}
 P(\I)=\I=Q(\I).
\end{equation}
Indeed, $P(\I)$ is a nonempty compact interval.  Moreover,
\[
 P(P(\I))\subseteq P(\I),\qquad
 Q(P(\I))=P(Q(\I))\subseteq P(\I),
\]
so $P(\I)$ is invariant under both maps.  Minimality forces $P(\I)=\I$,
and the proof for $Q$ is identical.

Put
\[
 H=P\circ Q=Q\circ P.
\]
The map $H$ is again continuous and finite-piece piecewise monotone.  By
Theorem~\ref{pm:thm:BF},
\begin{equation}\label{pm:eq:positive-entropy}
 h_{\rm top}(H)>0.
\end{equation}
Set
\[
 \lambda=\exp(h_{\rm top}(H))>1.
\]

\subsection{The branch-transfer operator with plateau intervals}

Let $J$ be a compact interval and let $u:J\to\I$ be continuous and
monotone.  Define the set of plateau values
\[
 \Pi(u)=\{z\in\I:u^{-1}(z)\text{ is nondegenerate}\}.
\]

\begin{lemma}[Image measures on a monotone branch]
\label{pm:lem:monotone-image-measure}
The set $\Pi(u)$ is countable.  If $\sigma$ is a finite signed nonatomic
Borel measure on $\I$, then:
\begin{enumerate}
\item $u(E)$ is Borel for every Borel set $E\subseteq J$;
\item the rule
\[
 \beta_{u,\sigma}(E)=\sigma(u(E))
\]
defines a finite signed nonatomic Borel measure on $J$.
\end{enumerate}
\end{lemma}

\begin{proof}
The interiors of the nondegenerate fibers of $u$ are pairwise disjoint
nonempty open intervals.  Assigning to each such interior a rational point
shows that $\Pi(u)$ is countable.

Put
\[
 J_0=J\setminus u^{-1}(\Pi(u)).
\]
The restriction $u_0=u|_{J_0}$ is an injective Borel map between standard
Borel spaces.  By the Lusin--Souslin theorem, $u_0(E\cap J_0)$ is Borel
whenever $E$ is Borel, and $u_0^{-1}$ is Borel on its image.  Moreover
\[
 u(E)=u_0(E\cap J_0)\cup\bigl(u(E)\cap\Pi(u)\bigr).
\]
The second set is countable, so $u(E)$ is Borel.

The nonatomic measure $\sigma$ vanishes on the countable set $\Pi(u)$.
Therefore
\[
 \sigma(u(E))=\sigma(u_0(E\cap J_0)).
\]
The right-hand side is the pullback, through the Borel bijection $u_0$, of
the restriction of $\sigma$ to $u_0(J_0)$, extended by zero across
$J\setminus J_0$.  It is consequently a finite signed countably additive
Borel measure.  Finally, the image of a singleton is a singleton, so this
measure is nonatomic.
\end{proof}

Let $T:\I\to\I$ be a continuous finite-piece piecewise-monotone map.
Choose closed intervals on which $T$ is monotone and whose union is $\I$.
Assign each common internal endpoint to exactly one of its adjacent
intervals.  The resulting pairwise disjoint Borel sets form a finite
partition $\Lap(T)$; we call them the \emph{assigned cells}.  Each assigned
cell differs from its underlying closed monotonicity interval by at most its
endpoints.

Let $\M(\I)$ denote the real vector space of finite signed nonatomic Borel
measures on $\I$.  Define
\begin{equation}\label{pm:eq:transfer-definition}
 (\T_T\sigma)(E)
 =\sum_{L\in\Lap(T)}\sigma\bigl(T(E\cap L)\bigr)
\end{equation}
for every Borel set $E\subseteq\I$.  Lemma~\ref{pm:lem:monotone-image-measure}
shows that each summand is a finite signed measure supported on its assigned
cell.

Refining a monotonicity partition does not alter
\eqref{pm:eq:transfer-definition}.  Indeed, if a monotone interval is divided
at one point, the images of its two parts overlap in at most one value; a
nonatomic measure ignores that overlap.  Hence $\T_T$ depends only on $T$.
It is real-linear on $\M(\I)$, preserves nonnegative
measures, and preserves nonatomicity.  If a Borel set $E$ is contained in
one assigned cell $L$, then only that cell contributes and
\[
 (\T_T\sigma)(E)=\sigma(T(E)).
\]
In particular, if $K$ is an interval contained in the underlying closed
monotonicity interval, endpoint assignments change both sides only by
nonatomic point masses, so
\begin{equation}\label{pm:eq:single-lap-scaling}
 (\T_T\sigma)(K)=\sigma(T(K)).
\end{equation}

\begin{lemma}[Composition law]\label{pm:lem:composition}
For continuous finite-piece piecewise-monotone maps $U,T$,
\begin{equation}\label{pm:eq:composition-law}
 \T_{U\circ T}=\T_T\circ\T_U,
 \qquad
 (\T_T\circ\T_U)(\sigma)=\T_T(\T_U\sigma).
\end{equation}
Consequently, if $U$ and $T$ commute, then $\T_U$ and $\T_T$ commute.
\end{lemma}

\begin{proof}
Use assigned cells for both maps.  For $L\in\Lap(T)$ and $M\in\Lap(U)$, the
nonempty sets
\[
 K_{L,M}=L\cap T^{-1}(M)
\]
form a finite pairwise disjoint Borel partition of $\I$.  Because $T$ is
monotone on the closure of $L$, each such set is an interval with some
endpoints assigned or deleted, and $U\circ T$ is monotone on its closure.
Thus, after deleting empty sets, these are the assigned cells of a valid
refinement for the composite.

For a Borel set $E$, using the definition and then this refinement gives
\begin{align*}
 ((\T_T\circ\T_U)\sigma)(E)
 &=\sum_{L\in\Lap(T)}(\T_U\sigma)(T(E\cap L))\\
 &=\sum_{L\in\Lap(T)}\sum_{M\in\Lap(U)}
   \sigma\bigl(U(T(E\cap L)\cap M)\bigr)\\
 &=\sum_{L,M}
   \sigma\bigl(U(T(E\cap L\cap T^{-1}(M)))\bigr).
\end{align*}
This is exactly $\T_{U\circ T}\sigma(E)$ for the refined monotonicity
partition; the middle equality uses the exact set identity
\[
 T(E\cap L)\cap M=T(E\cap L\cap T^{-1}(M)).
\]
If $U\circ T=T\circ U$, applying the formula in both orders gives
$\T_T\circ\T_U=\T_U\circ\T_T$.
\end{proof}

We now obtain the expanding eigenmeasure for $H$ from
Theorem~\ref{pm:thm:constant-slope}.  Let $\psi$ be the semiconjugacy there.
Because $\psi$ is continuous, nondecreasing, and onto,
$\psi(0)=0$ and $\psi(1)=1$.  The Lebesgue--Stieltjes premeasure
\[
 \mu((a,b])=\psi(b)-\psi(a),\qquad \mu(\{0\})=0,
\]
extends uniquely to a Borel probability measure on $\I$.  Continuity of
$\psi$ makes this measure nonatomic, and it satisfies
\begin{equation}\label{pm:eq:mu-from-psi}
 \mu([x,y])=\psi(y)-\psi(x)\qquad(0\le x\le y\le1).
\end{equation}
For every interval $K=[x,y]$ contained in one monotonicity piece of $H$,
\eqref{pm:eq:semiconjugacy-scaling} gives
\[
 \mu(H(K))=\lambda\mu(K).
\]
On a fixed closed monotonicity piece, Lemma~\ref{pm:lem:monotone-image-measure}
shows that $E\mapsto\mu(H(E))$ is a finite Borel measure.  It and
$E\mapsto\lambda\mu(E)$ agree on all closed subintervals by the preceding
display.  Those intervals form a generating $\pi$-system, so uniqueness of
finite Borel extension makes the two measures equal on the whole piece.
Using assigned cells, whose endpoint choices are invisible to the
nonatomic measures, and summing over the pieces yields
\begin{equation}\label{pm:eq:H-eigenmeasure}
 \T_H\mu=\lambda\mu.
\end{equation}

\subsection{Finite-dimensionality of the expanding eigenspace}

Define
\[
 \mathscr E_\lambda(H)
 =\{\sigma\in\M(\I):\T_H\sigma=\lambda\sigma\}.
\]
This is a real vector space of signed measures.

\begin{lemma}[Finite-dimensional expanding eigenspace]
\label{pm:lem:finite-dimensional}
The vector space $\mathscr E_\lambda(H)$ is finite-dimensional.
\end{lemma}

\begin{proof}
Choose $N\geq1$ with
\[
 \lambda^N>2.
\]
By Lemma~\ref{pm:lem:composition},
\[
 \T_{H^N}=\T_H^N.
\]
Let
\[
 0=c_0<c_1<\cdots<c_r=1
\]
contain all endpoints of a finite monotonicity partition for $H^N$.

For $\sigma\in\mathscr E_\lambda(H)$ put
\[
 S_\sigma(x)=\sigma([0,x])
\]
and define the linear map
\[
 \Xi:\mathscr E_\lambda(H)\longrightarrow\mathbb R^{r+1},
 \qquad
 \Xi(\sigma)=\bigl(S_\sigma(c_0),\ldots,S_\sigma(c_r)\bigr).
\]
We show that $\Xi$ is injective.

Suppose $\Xi(\sigma)=0$, and put
\[
 M=\sup_{x\in\I}|S_\sigma(x)|<\infty.
\]
Take $x\in[c_{i-1},c_i]$.  Since $H^N$ is monotone on this interval,
\eqref{pm:eq:single-lap-scaling} and
$\T_{H^N}\sigma=\lambda^N\sigma$ give
\[
 \lambda^N\sigma([c_{i-1},x])
 =\sigma\bigl(H^N([c_{i-1},x])\bigr).
\]
The left-hand interval begins at a point where $S_\sigma$ vanishes, so
\[
 \sigma([c_{i-1},x])=S_\sigma(x).
\]
The image on the right is an interval, possibly a singleton.  Since
$\sigma$ is nonatomic, the measure of any interval is the difference of two
values of $S_\sigma$; hence its absolute value is at most $2M$.  Therefore
\[
 |S_\sigma(x)|\leq \frac{2}{\lambda^N}M.
\]
Taking the supremum over $x$ yields
\[
 M\leq\frac{2}{\lambda^N}M.
\]
Since $2/\lambda^N<1$, we get $M=0$.  Thus $S_\sigma\equiv0$ and
$\sigma=0$.  Hence $\Xi$ is injective, proving
\[
 \dim\mathscr E_\lambda(H)\leq r+1<\infty.
\]
\end{proof}

For the rest of the reduction, equip $\mathscr E_\lambda(H)$ with the
coordinate norm
\begin{equation}\label{pm:eq:Xi-norm}
 \|\sigma\|_\Xi
 =\max_{0\leq i\leq r}|S_\sigma(c_i)|.
\end{equation}
Injectivity of $\Xi$ makes this a norm.  It gives the usual topology on the
finite-dimensional space $\mathscr E_\lambda(H)$; in particular every
linear functional and every linear endomorphism of this space is
continuous.

\subsection{A simultaneous conformal measure}

Write
\[
 A=\T_P,\qquad B=\T_Q.
\]
Since $P,Q,H$ commute pairwise, Lemma~\ref{pm:lem:composition} shows that
$A$ and $B$ commute with $\T_H$ and hence preserve
$\mathscr E_\lambda(H)$.  With operator composition written explicitly,
\begin{equation}\label{pm:eq:AB-lambda}
 B\circ A=A\circ B=\T_H=\lambda\,\id
 \quad\text{on }\mathscr E_\lambda(H).
\end{equation}
Indeed, $\T_H=\T_{P\circ Q}=\T_Q\circ\T_P=B\circ A$ by
\eqref{pm:eq:composition-law}; commutation gives the other order.

Let
\[
 \mathfrak P
 =\{\nu\in\mathscr E_\lambda(H):
       \nu\geq0,\ \nu(\I)=1\}.
\]
This set is nonempty because it contains $\mu$ from
\eqref{pm:eq:H-eigenmeasure}.  It is convex.

\begin{lemma}\label{pm:lem:compact-slice}
The set $\mathfrak P$ is compact in the finite-dimensional vector space
$\mathscr E_\lambda(H)$.
\end{lemma}

\begin{proof}
Use the injective coordinate map $\Xi$ from the proof of
Lemma~\ref{pm:lem:finite-dimensional}.  For every $\nu\in\mathfrak P$,
\[
 0\leq S_\nu(c_i)\leq1,
\]
so $\|\nu\|_\Xi\leq1$ and $\mathfrak P$ is bounded.  For each Borel set
$E$, evaluation $\nu\mapsto\nu(E)$ is a linear functional on the
finite-dimensional normed space and is therefore continuous.  The positive
cone is the closed intersection
\[
 \bigcap_{E\ {\rm Borel}}\{\nu:\nu(E)\geq0\},
\]
and the mass-one hyperplane is closed for the same reason.  Thus
$\mathfrak P$ is closed and bounded in a finite-dimensional space, and
hence compact.
\end{proof}

For $\nu\in\mathfrak P$, the measure $A\nu$ is nonnegative and belongs to
$\mathscr E_\lambda(H)$.  It is nonzero.  More explicitly, because $P$ is
onto and the images of its finitely many monotonicity pieces cover $\I$,
\begin{equation}\label{pm:eq:mass-lower-bound}
 (A\nu)(\I)
 =\sum_{L\in\Lap(P)}\nu(P(L))
 \geq \nu\left(\bigcup_{L\in\Lap(P)}P(L)\right)
 =1.
\end{equation}
Define
\[
 \Phi:\mathfrak P\longrightarrow\mathfrak P,
 \qquad
 \Phi(\nu)=\frac{A\nu}{(A\nu)(\I)}.
\]
The operator $A$ is a continuous linear endomorphism of
$\mathscr E_\lambda(H)$, the mass functional is continuous, and
\eqref{pm:eq:mass-lower-bound} keeps the denominator at least $1$.
Therefore $\Phi$ is a continuous selfmap of the nonempty compact convex set
$\mathfrak P$.  Brouwer's fixed-point theorem supplies
$\nu\in\mathfrak P$ with $\Phi(\nu)=\nu$.  Hence, for
\[
 C=(A\nu)(\I)>0,
\]
we have
\begin{equation}\label{pm:eq:P-eigenmeasure}
 \T_P\nu=C\nu.
\end{equation}
Using $B\circ A=\lambda\id$ on $\mathscr E_\lambda(H)$, we obtain
\[
 C\,B\nu=(B\circ A)\nu=\lambda\nu.
\]
Thus, with $D=\lambda/C>0$,
\begin{equation}\label{pm:eq:Q-eigenmeasure}
 \T_Q\nu=D\nu.
\end{equation}
We have proved the following central statement.

\begin{proposition}[Simultaneous conformal measure]
\label{pm:prop:simultaneous-measure}
There is a nonatomic Borel probability measure $\nu$ and constants
$C,D>0$ such that
\[
 \T_P\nu=C\nu,\qquad
 \T_Q\nu=D\nu,\qquad
 CD=\lambda>1.
\]
\end{proposition}

\subsection{The common monotone quotient}

Define the distribution function
\begin{equation}\label{pm:eq:theta}
 \theta(x)=\nu([0,x]).
\end{equation}
Because $\nu$ is nonatomic, $\theta$ is continuous.  It is nondecreasing,
and since $\nu$ is a probability measure,
\[
 \theta(0)=0,\qquad \theta(1)=1.
\]
Thus $\theta:\I\to\I$ is onto.

\begin{lemma}[Respect of fibers]\label{pm:lem:fibers}
The maps $P$ and $Q$ respect the fibers of $\theta$.  Consequently there
are unique continuous maps $\overline P,\overline Q:\I\to\I$ satisfying
\begin{equation}\label{pm:eq:factor-identities}
 \theta P=\overline P\theta,
 \qquad
 \theta Q=\overline Q\theta.
\end{equation}
The descended maps commute and are onto.
\end{lemma}

\begin{proof}
Suppose $x<y$ and $\theta(x)=\theta(y)$.  Then
\[
 \nu([x,y])=0.
\]
Partition $[x,y]$ at the finitely many monotonicity-partition endpoints of
$P$, obtaining closed intervals $J_1,\ldots,J_s$, each contained in the
underlying closed interval of one assigned cell.  Their only overlaps are
endpoints, which have zero measure.  Since $\nu$ is nonnegative and each
$J_i\subseteq[x,y]$, every $J_i$ has measure zero.  By
\eqref{pm:eq:single-lap-scaling} and
\eqref{pm:eq:P-eigenmeasure},
\[
 \nu(P(J_i))=C\nu(J_i)=0.
\]
The continuous image $P([x,y])$ is the finite union of the intervals
$P(J_i)$, so it has $\nu$-measure zero.  In particular, the interval with
endpoints $P(x)$ and $P(y)$ has measure zero, and therefore
\[
 \theta(P(x))=\theta(P(y)).
\]
Thus $P$ respects the fibers of $\theta$.  The proof for $Q$ is identical,
using \eqref{pm:eq:Q-eigenmeasure}.

Since a continuous surjection from a compact space to a Hausdorff space is
a quotient map, the well-defined formulas
\[
 \overline P(\theta(x))=\theta(P(x)),\qquad
 \overline Q(\theta(x))=\theta(Q(x))
\]
give the unique continuous factor maps in
\eqref{pm:eq:factor-identities}.  Their commutation follows from
\[
 \overline P\overline Q\theta=\theta PQ
 =\theta QP=\overline Q\overline P\theta
\]
and onto-ness of $\theta$.  Finally, given $t\in\I$, choose $x$ with
$\theta(x)=t$ and then $z$ with $P(z)=x$.  The factor identity gives
$t=\overline P(\theta(z))$, so $\overline P$ is onto; the proof for
$\overline Q$ is identical.
\end{proof}

\begin{lemma}[Piecewise-linear quotient]\label{pm:lem:PL-quotient}
The map $\overline P$ is a continuous finite-piece piecewise-linear map
whose nonzero slopes have absolute value $C$.  Likewise, $\overline Q$ is
finite-piece piecewise linear with nonzero slopes of absolute value $D$.
Every plateau interval of either original map is collapsed by $\theta$.
\end{lemma}

\begin{proof}
Let $L=[a,b]$ be a closed monotonicity piece of $P$.  Its image under the
nondecreasing map $\theta$ is the interval
\[
 \theta(L)=[\theta(a),\theta(b)],
\]
possibly degenerate.  Suppose it is nondegenerate, and take
$s<t$ in $\theta(L)$.  Choose $x<y$ in $L$ with
\[
 \theta(x)=s,\qquad \theta(y)=t.
\]
By \eqref{pm:eq:single-lap-scaling} and \eqref{pm:eq:P-eigenmeasure},
\[
 \nu(P([x,y]))=C\nu([x,y])=C(t-s).
\]
If $P$ is nondecreasing on $L$, then $P([x,y])=[P(x),P(y)]$ and
\[
 \overline P(t)-\overline P(s)
 =\theta(P(y))-\theta(P(x))
 =C(t-s).
\]
If $P$ is nonincreasing, the same computation gives
\[
 \overline P(t)-\overline P(s)=-C(t-s).
\]
Thus $\overline P$ is affine with slope $+C$ or $-C$ on every
nondegenerate interval $\theta(L)$.  The images of the finitely many
consecutive monotonicity pieces form a finite interval partition of $\I$,
after degenerate members are removed.  Hence $\overline P$ is finite-piece
piecewise linear.

Let $J$ now be an arbitrary plateau interval of $P$, without assuming that
it lies in one chosen monotonicity piece.  Subdivide $J$ at all
monotonicity-partition endpoints.  Each resulting closed subinterval $K$
lies in one underlying assigned cell, and
\[
 C\nu(K)=\nu(P(K))=0.
\]
The subdivision is finite and its overlaps are nonatomic endpoints, so
$\nu(J)=0$ and $\theta$ is constant on $J$.  Thus every plateau interval is
collapsed.  The proof for $\overline Q$, including the affine-slope and
plateau arguments, is the same.
\end{proof}

\subsection{Conclusion}

By Lemmas~\ref{pm:lem:fibers} and \ref{pm:lem:PL-quotient}, the maps
$\overline P,\overline Q$ are commuting continuous finite-piece
piecewise-linear selfmaps of $\I$.  Proposition~\ref{prop:PL-core} gives a point
$q\in\I$ such that
\[
 \overline P(q)=q=\overline Q(q).
\]
The fiber
\[
 K=\theta^{-1}(q)
\]
is a nonempty compact interval.  The factor identities imply
\[
 P(K)\subseteq K,
 \qquad
 Q(K)\subseteq K.
\]
By the minimality of the common invariant core chosen above, we must have
$K=\I$.  Then $\theta$ is constant, contradicting the fact that
$\theta:\I\to\I$ is onto.  This contradiction proves
Theorem~\ref{thm:main}.\hfill\(\square\)\par

\begin{thebibliography}{99}

\bibitem{Boyce}
W.~M. Boyce,
\emph{Commuting functions with no common fixed point},
Trans. Amer. Math. Soc. \textbf{137} (1969), 77--92,
doi: \href{https://doi.org/10.1090/S0002-9947-1969-0236331-5}
{10.1090/S0002-9947-1969-0236331-5}.

\bibitem{Huneke}
J.~P. Huneke,
\emph{On common fixed points of commuting continuous functions on an
interval},
Trans. Amer. Math. Soc. \textbf{139} (1969), 371--381,
doi: \href{https://doi.org/10.1090/S0002-9947-1969-0237724-2}
{10.1090/S0002-9947-1969-0237724-2}.

\bibitem{Navas}
A.~Navas,
\emph{Groups of Circle Diffeomorphisms},
Chicago Lectures in Mathematics,
University of Chicago Press, Chicago, 2011;
Propositions 2.2.2 and 2.2.3 and Theorems 3.1.1 and 4.1.1,
\url{https://arxiv.org/abs/math/0607481}.

\bibitem{Shultz}
F.~Shultz,
\emph{Dimension groups for interval maps II: the transitive case},
Ergodic Theory Dynam. Systems \textbf{27} (2007), 1287--1321;
Propositions 3.6 and 4.3 and Corollaries 4.4 and 4.6,
\url{https://arxiv.org/abs/math/0405467}.

\bibitem{Parry}
W.~Parry,
\emph{Symbolic dynamics and transformations of the unit interval},
Trans. Amer. Math. Soc. \textbf{122} (1966), 368--378,
doi: \href{https://doi.org/10.1090/S0002-9947-1966-0197683-5}
{10.1090/S0002-9947-1966-0197683-5}.

\bibitem{Durand}
F.~Durand,
\emph{Cobham's theorem for substitutions},
J. Eur. Math. Soc. \textbf{13} (2011), 1799--1814,
doi: \href{https://doi.org/10.4171/JEMS/294}{10.4171/JEMS/294},
\url{https://ems.press/journals/jems/articles/4506}.

\bibitem{FugateMcLean}
J.~B. Fugate and T.~B. McLean,
\emph{Compact groups of homeomorphisms on tree-like continua},
Trans. Amer. Math. Soc. \textbf{267} (1981), 609--620,
doi: \href{https://doi.org/10.1090/S0002-9947-1981-0626493-5}
{10.1090/S0002-9947-1981-0626493-5}.

\bibitem{AlsedaMisiurewicz}
Ll.~Alsed\`a and M.~Misiurewicz,
\emph{Semiconjugacy to a map of a constant slope},
Discrete Contin. Dyn. Syst. Ser. B \textbf{20} (2015), no.~10,
3403--3413,
doi: \href{https://doi.org/10.3934/dcdsb.2015.20.3403}
{10.3934/dcdsb.2015.20.3403}.

\bibitem{BogatyiFrolkina}
S.~A. Bogatyi and O.~D. Frolkina,
\emph{A common fixed point of commuting mappings of a tree},
Vestnik Moskov. Univ. Ser.~1 Mat. Mekh. 2002, no.~6, 3--10;
English translation in Moscow Univ. Math. Bull. \textbf{57} (2002),
no.~6, 1--8 (2003),
\url{https://www.mathnet.ru/vmumm1423}.

\end{thebibliography}

\end{document}
